% =====================================================================
%  Harald Høffding, On Some Recent Works in the Philosophy of Religion (1903)
%  Translated from the Danish: Om nogle religionsfilosofiske Arbejder fra den nyeste Tid.  Indbydelsesskrift til
%  Kjøbenhavns Universitets Aarsfest i Anledning af Hans Majestæt Kongens Fødselsdag den 8. April 1903.
%  Kjøbenhavn: Universitetsbogtrykkeriet (J. H. Schultz), 1903.
%
%  TRANSLATION -- GENERATED; DO NOT EDIT BY HAND.  Built by
%    python3 ebook-method/religionsfilosofiske-arbejder-1903/tr1903.py build --write
%  from this template, the English of the 1905 reprint (ebook-method/mindre-arbejder-2/tr/ARB-A..D.texfrag), en_edits.py,
%  en_notes.py and anchors-en.txt.  Made from the collated transcription texts/hoeffding/religionsfilosofiske-arbejder-1903/
%  transcription.tex; [N] in the margin = the page of the 1903 print.
%
%  STATUS: COMPLETE (2026-10-26).
%
%  METHOD (Hans Halvorson, 2026-10-26).  This print was reprinted, revised by Høffding, in Mindre Arbejder. Anden Række
%  (1905), pp. 58-140, whose English already exists (Shorter Works. Second Series; ebook-method/mindre-arbejder-2/tr/
%  ARB-A..D.texfrag, translated 2026-10-16 with ebook-method/mindre-arbejder-2/GLOSSARY.md).  That English is REUSED
%  wherever 1903 and 1905 agree, and changed or newly translated only where the 1903 Danish differs (tr1903.py work lists
%  every 1903 paragraph whose words differ from 1905; ebook-method/religionsfilosofiske-arbejder-1903/variants.tsv lists
%  the differences).  Paragraphs, headings, italics and page markers follow the 1903 Danish one to one; the notes are
%  endnotes as in 1903.  Same glossary as the Second Series.  [N] in the margin: the 1903 page (anchors-en.txt, placed
%  at the English for the first full Danish word of each page).
%  NEWLY TRANSLATED (not in the 1905 English): the Notes (all 13; pp. 84-86), and in the text: p. 4 "whose standpoint is
%  otherwise very different from mine" (1905 names the critic instead); p. 9 "new creations" (Nyskabelser; 1905
%  Nydannelser); p. 15 "for the sufficiency and vitality of the objective being of religion"; p. 21 "p. 173" (1905
%  "§ 66"); p. 23 "a mere phenomenon"; p. 38 "its extreme forms"; p. 42 "and the formula was created"; p. 57 "there was
%  perhaps"; p. 63 "sees the doctrine as fixed"; p. 71 "in construction, in
%  comparison or in criticism".  Removed (1905 only): the parenthesis "(DOWDEN: ...)" on p. 10 (in 1903 note 4), the 1905
%  footnote (Moderne Filosofer).  Structural: the paragraph on Goethe's Faust (p. 33/34) stands apart as in 1903; the
%  paragraph broken at the 1905 page turn 123 ("Already the Greek | philosophy's") is one paragraph (p. 65).  The 1903
%  endnote marks (1-13) set where the Danish has them.  Italics brought to the 1903 print in 22 places (en_edits.py).
%  NOT CHANGED: differences of punctuation, word division and spelling that do not alter the English; where 1903 has
%  commas and 1905 stops (p. 46 "(revivals), Nogle ...") the English keeps its sentences.
%  The 1905 English had rendered by the sense several places where the 1905 print or its transcription is faulty (p. 64
%  "her" for the misprint "hor", p. 59 "intuitions"); these agree with 1903 and are kept.  Two faults of the 1905 English
%  found here were mended in ARB-D itself (2026-10-26): "one could then define religion" and the paragraph at its p. 123.
%  Every edit is in ebook-method/religionsfilosofiske-arbejder-1903/en_edits.py with its reason; the notes in en_notes.py.
%  Rebuild: python3 ebook-method/religionsfilosofiske-arbejder-1903/tr1903.py build --write (after the Danish).
% =====================================================================
\documentclass[12pt,a4paper,openany]{book}
\usepackage[utf8]{inputenc}
\usepackage[T1]{fontenc}
\usepackage{libertinus}
\usepackage{libertinust1math}
\usepackage[german,french,english]{babel}
\usepackage[protrusion=true,expansion=true]{microtype}
\usepackage{setspace}
\usepackage[margin=1.3in]{geometry}
\usepackage{fancyhdr}
\usepackage{marginnote}
\usepackage{hyperref}
\hypersetup{pdftitle={On Some Recent Works in the Philosophy of Religion}, pdfauthor={Harald Høffding}, hidelinks}
\pagestyle{fancy}
\fancyhf{}
\fancyhead[LE,RO]{\thepage}
\fancyhead[C]{\textit{H. Høffding: On Some Recent Works in the Philosophy of Religion (1903)}}
\fancypagestyle{plain}{\fancyhf{}\fancyhead[LE,RO]{\thepage}}
\setlength{\headheight}{15pt}
\setstretch{1.3}
\newcommand{\opage}[1]{\marginnote{\footnotesize\textit{[#1]}}}
\newcommand{\kapitel}[3]{\chapter*{\centering\normalsize\if\relax\detokenize{#1}\relax\else #1\\[1ex]\fi\Large #2}\addcontentsline{toc}{chapter}{#3}}
\newcommand{\afsnit}[2]{\section*{\centering\normalsize\bfseries #1}\addcontentsline{toc}{section}{#2}}
\newcommand{\underafsnit}[2]{\subsection*{\centering\normalsize\normalfont\itshape #1}\addcontentsline{toc}{subsection}{#2}}
\newcommand{\streg}{\begin{center}\rule{3em}{0.4pt}\end{center}}
\newcommand{\ix}[3]{\par\noindent\hspace*{#1}#2\dotfill #3}
\pdfstringdefDisableCommands{\renewcommand{\opage}[1]{}\renewcommand{\textsuperscript}[1]{}}

\begin{document}

\frontmatter
\begin{titlepage}\centering
{\large INVITATION}\\[1ex]{\small TO}\\[1ex]{\large THE ANNUAL FESTIVAL OF THE UNIVERSITY OF COPENHAGEN}\\[1ex]
{\small ON THE OCCASION OF}\\[1ex]{\large HIS MAJESTY THE KING'S BIRTHDAY}\\[1ex]{\small 8 APRIL 1903}\par\vspace{2cm}
{\large HARALD HØFFDING}\\[1ex]{\large ON SOME RECENT WORKS IN THE PHILOSOPHY OF RELIGION}\par\vfill
{\small COPENHAGEN}\\[0.5ex]{\small PRINTED BY THE UNIVERSITY PRINTING OFFICE (J. H. SCHULTZ)}\\[0.5ex]1903\\[3em]
{\small [Om nogle religionsfilosofiske Arbejder fra den nyeste Tid. Translated from the Danish]}
\end{titlepage}
\chapter*{CONTENTS}
\noindent\hfill Page
\ix{0em}{INTRODUCTION}{1}
\ix{0em}{I. SYSTEMATIC ATTEMPTS}{12}
\ix{1.5em}{1. Hermann Siebeck}{12}
\ix{1.5em}{2. Auguste Sabatier}{18}
\ix{1.5em}{3. Rudolph Eucken}{21}
\ix{0em}{II. INVESTIGATIONS IN THE PSYCHOLOGY OF RELIGION}{32}
\ix{1.5em}{1. On Religious Judgments. (Max Reischle).}{35}
\ix{1.5em}{2. The Pathology of Religious Feeling. (Murisier).}{38}
\ix{1.5em}{3. Contributions to the Psychology of Confessional Religiosity.}{41}
\ix{3em}{a. Breakthrough and Continuity. (Starbuck).}{46}
\ix{3em}{b. The Influence of Temperament. (Coe).}{51}
\ix{3em}{c. Differences within One Confession. (Leuba, Arréat).}{54}
\ix{3em}{d. Poetry and Religion. (Santayana)}{65}
\ix{3em}{e. Comparison between Different Series of Types}{67}
\ix{1.5em}{4. Comparative Psychology of Religion. (James).}{69}
\ix{3em}{a. The Three Points of View}{71}
\ix{3em}{b. The Nature of Religious Experience}{75}
\ix{3em}{c. The Most Important Types}{76}
\ix{3em}{d. The Fundamental Religious Idea.}{79}
\ix{3em}{e. Remarks on James’s Philosophy of Religion}{80}
\par\medskip\noindent Notes p. 84.
\clearpage

\mainmatter
\kapitel{}{Introduction.}{Introduction}\opage{1}

By philosophy of religion I understand not an attempt to put speculation in the place of religion, still less an attempt to ground religious assumptions by the road of speculation, but an attempt to throw light on the position of religion in the intellectual life of the age. Religion is here the object of investigation as an actually given form of spiritual life, whose nature, significance and justification it is our task to determine. If one now asks what points of view philosophy commands for use in such an investigation, there can surely be talk of only three points of view, --- but these are really also all the points of view that philosophy commands at all: the logical (or epistemological), the psychological and the ethical\textsuperscript{1}). Since every religion --- whether one takes the word in a wider or a narrower sense --- finds expression in certain ideas, the question naturally arises of the validity of these ideas, their value as knowledge. Do they contain a consistent and valid knowledge of existence, and can they be brought into agreement with the ideas that science by its own roads has arrived at holding fast as valid? Even if such an \textit{epistemological} investigation should lead to the result that the significance of religious ideas cannot consist in procuring for us a scientific understanding of existence, the question of their value and justification is not thereby settled. It may rest on their relation to other sides of mental life than the purely intellectual. And \opage{2}direct observation also shows, indeed, that feeling, need and will play quite as great a role in religion as knowledge. There is therefore every occasion to undertake a \textit{psychological} investigation of the religious life, as observation teaches us to know it in the present and in historical tradition in the past. This investigation aims, first, at finding the constantly operative motives in religion and separating them from the changing forms and connections in which they appear under different conditions. Only in these motives can the essential in religion be sought, --- that which, so long as it possesses vitality, will always be able to express itself in new shapes. There exists no religion in general; but experience shows us certain fundamental motives in a wealth of different connections or variations. Therefore it is, second, the task of the psychology of religion to seek out and describe the most important types of religion and religious life. If, then, motives and types have thus been demonstrated, it is the business of the \textit{ethical} investigation to decide what value religion, as a peculiar form of human spiritual life, has for human life in its whole, great connection. The standard here is the same that we also apply otherwise when we discuss the value of a human quality or action. We examine their relation to certain fundamental values, in the last instance to a single fundamental value which for us is the decisive one. What this fundamental value is, it is the business of ethics to find out. If there is disagreement about the fundamental value that determines what we in the individual cases call valuable, there will also be disagreement about the valuation of religion. The question would of course be decided in advance if one assumed that religion (i.e. the religion one has or believes oneself to have) is itself the fundamental value by which everything is to be measured. But then there would likewise exist no religious problem, nor would there be any possibility of or need for a philosophy of religion, just as little as there would be any possibility of or need for a theory of knowledge, if one assumed that some religion or other contained the absolute truth, by which everything else one would call true was to be measured. A problem arises here, on the other hand, as soon as one becomes aware that in acknowledging religion \opage{3}one presupposes a certain standard as valid, a standard that is not drawn from religion alone, and especially when one discovers that in the religious domain one applies a standard that differs from that which one applies in other domains.

Of the three points of view from which the religious problem can and must be treated, it is the psychological that in the most recent times especially attracts attention. It is of course nothing new that religion is conceived as a psychological phenomenon. Every more profoundly reaching religious movement had to lead to it, especially by making it clear that what matters is not merely to stand in a relation of ideas to the Divine, but also to be moved by it in inward feeling. More and more it is pointed out from the religious side that the ultimate foundation of religion is human experience of life. The psychology of religion fully acknowledges this, and it then sets itself the task of describing this experience. It proceeds here from the view that, if the appeal to experience is to have any meaning, one must understand by experience immediate experience, and it then matters to determine what and how much can be immediately experienced. It is necessary to distinguish between what is really immediately given and what is added as explanation. The explanation is drawn from a circle of ideas that lies close at hand to man. And from the explanation the mode of expression must again be distinguished. But these three things, experience, mode of expression and explanation, it is extraordinarily difficult to keep apart from one another in the case of religious experience, which is essentially an experience of feeling. For the psychology of religion it therefore matters to seek out simple examples where the difference is clear, and with their help to explain to oneself the more complicated cases. But so much will be clear at once, that there is a long way from an immediate religious experience to the often very subtle dogmas or to the poetically shaped myths and legends that the historically given religions retain in their traditions. And a quite special distinction must be drawn between experiences that are evoked by such traditional ideas in the attempt to convert them into self-won experience, and experiences that at first hand lead to \opage{4}the formation of certain determinate ideas. Only this latter kind is original religious experience. And even with it one must still ask whether the one and the same experience could not be expressed and explained by other ideas and in other forms than is actually the case. Only if such a counter-test could be made, and if it then led to a negative result, would the proof from experience have been furnished.

I believe that this psychological method will acquire very great importance in the continuing religious discussion. It is with full right that appeal is made from the religious side to personal experience of life as the last criterion against the often very superficial appeal to the supposed results of reason or science. But if one does not treat the very concept of experience with strict criticism, distinguishing sharply between the given and its interpretation, one introduces a dogmatism that is far worse than the old rationalistic dogmatism on which one loftily looks down, as surely as it is worse to deceive oneself and others with regard to the most inward relation of one's life than to commit logical errors in an intricate piece of reasoning. In any case psychology cannot allow itself to be prevented from attempting to analyze the phenomena, immediate and uncompounded at first glance, that are appealed to. There lies here a great material ready to hand in the utterances of the religious consciousness at different times and under highly different cultural conditions, attractive and refreshing especially for those who have had enough of the hitherto current way of treating the religious questions, both among the assailants and among the defenders.

It is this that I have sought to the best of my ability to impress and carry through in my Philosophy of Religion, whose most extensive section is precisely psychological. I have also had the satisfaction that a theological critic, whose standpoint is otherwise very different from mine, has declared that my method is the only right one\textsuperscript{2}). But my astonishment must then be all the greater when I see that this same critic disputes my results --- by means of a quite different method than that which he has thus acknowledged! For he uses the concept of experience precisely in the vague and wide sense in which it is often used in everyday speech. He puts the question throughout thus: do certain traditional \opage{5}dogmas have value for life? But a consistently carried-out investigation in the psychology of religion goes also the opposite way and asks: do life values really disappear because one does not accept certain dogmas, or because one puts poetic symbols in the place of the dogmas? I myself dare (although I do not hide my own personal standpoint in this matter) not give a definite answer to this question; I have expressly emphasized the possibility that there might here be a loss of value. This possibility no science can refute, and therefore the assumption of the persistence of value rests on a faith, even if one does not extend this assumption beyond the limits of the human world. ---

In a domain where different tendencies still clash --- and perhaps for all eternity --- it will be of great importance to investigate whether the procedure employed is the only one that can come into question. Therefore I have wished to bring forward a series of attempts in the philosophy of religion, most of which first appeared after the completion of my book, in order to see whether there are other points of view to be adopted than the three mentioned (the logical, the psychological and the ethical), and specifically to see whether the attempted determination of religious experience and of the fundamental idea of religion should need supplementing or correcting.

Not only as regards the method but also as regards the results, such a comparison with the contemporary works of other investigators may be of importance. I therefore permit myself to recall that the results at which I believe I have arrived are mainly the following three. --- The religious problem falls under the general problem of valuation, and the method for its treatment will therefore more and more become ethical valuation alongside psychological description. Neither of these two procedures can be dispensed with, since one can of course only value accurately what is accurately described. --- The religious need may be characterized as a need to assert the persistence of value, and religious experience consists of what man has experienced concerning the relation between the values he knows and the reality he knows. The special religious assumptions indicate --- whether they appear in the form of myth, dogma, legend or symbol --- the ways in which \opage{6}men have sought partly to express, partly to explain what they have thus experienced. --- In the conditions of intellectual life in modern times there is nothing to prevent such a need from continuing to express itself and to find satisfaction, even if it must more and more happen in other forms than in earlier times, --- and even if it is not necessary that it express itself in everyone. It is a hypothesis of the history of religion that the proposition of the persistence of value is the fundamental religious idea, i.e. that it is this which expresses the pervading tendency in the historically given forms of religion. But the proposition of the persistence of value can also be set up as a philosophical hypothesis. It denotes a view of life that may claim to be \textit{compatible with} a strictly carried-through scientific world-view, even if it cannot be \textit{derived} from it. The view of life stands in everyone with a certain independence in relation to the world-view. There takes place a constant influence of the latter on the former; but never is the view of life a mere product of the scientifically grounded world-view. For the view of life contains a plus which derives from the personality's valuation of the existence it seeks to understand with the help of science. Now there may indeed be found within valuation greater or lesser consistency, and the points of view that are laid at the basis of valuation may be more or less natural and essential. But in the domain of valuation we are to a far higher degree compelled to be empiricists than where understanding and explanation are concerned. The irrational (in the mathematical sense), which our thinking must always leave standing over, has a greater extent in the domain of valuation than in that of pure theory. Beyond the \textit{possibility} of a view of life that builds on the fundamental idea of the persistence of value, philosophy cannot lead us\textsuperscript{3}). The fundamental motive of every view of life springs from practical experience of the fate of values in the world of reality; and if this experience leads to faith in the persistence of value, then this faith will not need to stand as a mystical hypothesis, a mere object of passive contemplation, but will naturally become a working hypothesis, in that it spurs one to find new values when the time of the old is past, or perhaps even oneself \opage{7}to produce such. It then gives courage to work both in the domain of thought and in that of reality in order to win progressive confirmation. Perhaps not all individuals will be led by their experiences to harbor such a faith; it may suffice if it maintains itself in the race. For the life of the race, in order to be rich and vigorous, needs many different individualities, perhaps even a constant struggle between spiritual tendencies. Many individuals may be so absorbed by the special work of thought and of reality that they do not have time, strength and interest left over to form any view of life expressed in determinate thoughts. Here, as in other respects, a division of labor may take place. Why should the division of labor be restricted to the purely external conditions of life? Is there not reason to assume that it intervenes in the most inward concerns of men? Life will become richer thereby, just as the world of art becomes richer the more the individual artist follows his individual need and capacity. The very proposition of the persistence of value perhaps even requires that not all believe in it or think of it.

These last remarks I lay particular weight on in face of an objection that has been directed against my book: that it presents a purely subjective conception. I have not hidden my own personal position toward the questions I discuss; but the main thing for me has been to investigate the possibilities that present themselves under the cultural conditions of the present in a religious respect, and here it was not enough to confine oneself to a single individual. It is a comparative investigation that I have wished to make. I have wished to undertake a comparison between the position of the religious in the past and in the present, and between its significance for different personalities and standpoints in the present. What matters in a philosophy of religion is not the philosopher's own faith, but the possibilities of faith that he finds in the intellectual life of the race, as it can now be carried on in clarity and in depth. Less and less has a philosopher the right to assume that the psychology of the race is exhausted with his own individual psychology.

A related objection has been raised against my presentation by \opage{8}an energetic German scholar, whom I have the honor to be able to call my personal friend. In a detailed review of my book in the journal ``Das freie Wort'' (1902) FERDINAND TÖNNIES has criticized my conception of the religious problem in a manner that seems to me capable of illuminating the relation between the history of religion and the philosophy of religion.

``For Guyau,'' he says, ``and I confess that in this I side with him, religion is in the first instance an object of social psychology or sociology. For Høffding it belongs almost entirely to individual psychology, which indeed coincides with what one is accustomed to understand by psychology at all. He calls Buddhism and Christianity `the higher popular religions': in them the inner psychological drama is to appear as the real world drama . . . Now it is indeed unmistakable that in both these highly important religious systems the individual need and the individual thinking have acquired an overwhelming significance. But whence does this struggle of the soul derive? Is it an original and essential constituent of religion as popular belief and popular custom, or at any rate of the religious feeling that `bears' belief and custom? \textit{Or} does it, precisely as regards its purer and higher content, issue from another source, namely from a thinking, a philosophizing, that by its nature is alien to popular belief and popular custom, . . . although it has certainly attained in those two popular religions an adjustment with popular and wild-growing (naturwüchsige) ideas and feelings, indeed in Christianity has even concluded an eternal league with them?''

The consequence of this line of thought seems bound to be that Buddhism and Christianity are not really religions. Indeed, one would really have to go further. Is there any positive religion whatever in which elements of thinking cannot be demonstrated? The most elementary, not to say the most primitive religions build on animism, man's most elementary philosophy, and thus at all stages of development a people's religion will bear the stamp of its antecedent ``philosophy''. Historians sometimes have rather romantic ideas of the immediacy and ``Naturwüchsigkeit'' with which men enter into their cult, myths and dogmas. \opage{9}There take place, as the history of religion itself shows, constant combinations, supplementations and metamorphoses, in which a more or less profound work of reflection co-operates. But what for me is the main point: it is misleading here to speak merely of thinking and philosophy. It is first and foremost personal experience of life that exerts its influence in these changes, which are often new creations. The personalities who have been decisive for the origin and development of the higher popular religions have first and foremost laid their own experiences of the valuable and its fate in the world at the basis of the changes they were led to make in what was handed down. No religion makes itself. Religion arises and exists as religion only in the minds and thought of men. It must have been certain determinate human beings who shaped cult, myth and dogma, just as it must have been certain determinate human beings who composed the ``popular poetry'', and just as ``society'' at all times consists of individual human beings. To that extent all psychology must be individualistic, even if it zealously applies the sociological method (alongside other psychological methods), which there is quite special occasion to do in the psychology of religion. --- At the end of his review Tönnies, moreover, after having repeated that in religion he sees especially an object for social psychology, says: ``Ich bekenne meinerseits, dass mich Høffding gelehrt hat, die Tragweite der individuellen Entscheidung besser als bisher zu würdigen''. ---

Only when one lays weight on the subjective element in religion does it become possible to find an inner continuity between the changing religions and views of life. And on this element rests also the significance for the continued life of the race that a religion can have, even though its time is once past. It is the great human treasures that the religious spirit has been able to raise in the inner life of men that constitute the most precious inheritance that the religions of the past leave behind. A poem like BUNYAN'S ``Pilgrims Progress'', which has been called the divina comedia of Protestantism, has thus acquired a lasting significance (at least for the English race) beyond the limits of the definite confession within which it came into being. And the reason is that it \opage{10}expresses what moved in the soul of a human being deeply gripped by one of life's oppositions. An excellent historian of literature has rightly said of it: ``All that is best and most characteristic in Bunyan's writings proceeds from the inner drama in which the actors were God, Satan and a lonely human soul. Institutions, churches, ordinances, rites and ceremonies help him little or not at all . . . . Let everyone faithfully hold fast to his most inward experiences in the great concerns of life, and his words will awaken manifold echoes in other souls; \textit{the deepest fellowship is found not in institutions or corporations or churches, but in the lonely heart's secret}.''\textsuperscript{4}). --- In my Philosophy of Religion I have also pointed out how the sense of the significance of the purely subjective elements in the religious domain can pave the way for a deeper feeling of community than that which is conditioned by union around common dogmas. To one and the same dogma or one and the same cult there may correspond highly different personal moods and states. This diversity is a proof of the independence of the subjective element. The relation between the objective and the subjective elements (or perhaps better: between the social and the individual elements) in religion is constantly varying, so long as there is life in religion. In the historically shaped religions it is the objective, social elements that most strongly strike the eye. And if the expression ``positive religion'' is to be understood in analogy with expressions like ``positive law'', what is meant thereby is a religion that appears as an institution, with a shaped cult and a fixed tradition. Psychologically viewed, however, what motivates man's adherence to and holding fast to what has thus been fixed, and what under certain conditions leads to making changes therein, must always be the essential thing. The significance of this subjective element can rise so that the objective element is as good as eliminated. We then come to a limit where there can be dispute whether the concept of religion can still be applied. It may easily become a dispute about words. It may lead to raising profound questions like that which a man put: ``Is the star that the astronomers call Uranus really Uranus?'' --- If one lays the \opage{11}weight on the fundamental motives, not on the external marks, it will be justified to extend the concept of religion beyond the sphere of the so-called positive religions. Religion then becomes a generic concept, within which religion as positive religion, built on tradition and provided with fixed forms of cult, denotes only a particular species. Thereby the oppositions are not underestimated; for within one and the same generic concept the strongest oppositions can appear, just as in reality a struggle of life and death can be waged between species that belong to the same genus. ---

After these introductory remarks I now go on to discuss some works in the philosophy of religion from the most recent times in order thereby to give a critical supplement to my own attempt. I shall choose such works as can be said both to be typical of the philosophy of religion of our time and to be instructive for those who are interested in the religious problem. --- There are, however, some authors whom it might in and of itself be natural to bring forward here, but whom I would rather discuss on another occasion, because they will best be understood in connection with the most recent philosophy in general. To these belong especially Guyau and Nietzsche, whom I shall in no long time have occasion to discuss in a work on modern philosophers, on which I am at work, and to which the present treatise can also be regarded as a supplement.

In the first section of the treatise I deal with some systematic attempts; in its second section I bring forward a series of more specialized investigations in the psychology of religion, especially from American literature.

\kapitel{I.}{Systematic Attempts.}{I. Systematic Attempts}\opage{12}

\afsnit{1. Hermann Siebeck.}{1. Hermann Siebeck}

SIEBECK (b. 1842), who works as Professor of Philosophy in Giessen, has made himself known by an excellent work on the history of psychology in antiquity and the Middle Ages and by his aesthetic studies. His chief work is the Philosophy of Religion that appeared in the year 1893. Originally proceeding from Herbart's school, Siebeck has later received influences from Lotze; but above all the effect of zealous Kant studies is felt in his chief work.

The task of the philosophy of religion is for Siebeck posed by the relation between religion and culture. Modern culture is inclined to regard itself as an end; religion it can, at best, acknowledge only as one of its imperfect predecessors. From the religious side, on the other hand, it is maintained that the life of culture --- like life and work in the world of experience in general --- can only be a means for a higher, ``supramundane'' end. The dispute that thus arises it is the business of the philosophy of religion to decide. It is therefore not a religiously attuned philosophy, like Neoplatonism in antiquity or like Schelling's later doctrine; nor is it, like the scholasticism of the Middle Ages, a presentation of a certain determinate religion in philosophical form. It regards religion in its various shapes as a fact in human life and investigates what significance this fact has for the development of human life. It seeks to determine in what domain of human spiritual life religion particularly belongs, and what \opage{13}influence it exerts from there on other sides of life. Philosophy here investigates religion in a similar manner to that in which in its other disciplines it investigates science, art, law and morality.

According to Siebeck it must essentially be an indirect method that the philosophy of religion can employ. It must test whether the very existence and development of the life of culture do not become worthless and self-contradictory when it cuts itself off from the supplementary points of view and motives that religion affords. A positive determination of the higher world to which religion will open access, philosophy cannot give. Further than to demonstrate the necessity of a supplementation it cannot reach; but it can test whether the forms in which the religious consciousness has sought to express the Divine are suited to furnish the supplementation that a critique of cultural development leads one to demand. In advance Siebeck declares as his conviction that in the indicated sense of the word a demonstration can be given of ``the normality of the essential content of the religious consciousness''.

For the philosophy of religion, when its task is determined in this way, it will be natural to keep to the highest forms of religion, since only they become of significance for the relation between religion and culture in modern intellectual life. Siebeck distinguishes between three main forms of religion: natural religion, ethical religion (``Moralitätsreligion'') and religion of redemption. In the first form the gods are conceived as powers of nature, and the relation to them is determined essentially by the need of self-preservation. In the second form the gods have a calling to fulfill, and their worship is motivated by the dignity this calling gives them. The third form, the religion of redemption, which appears typically in Buddhism and in Christianity, is characterized by the significance that prophetic personalities acquire. Myth and cult, which rule in the earlier forms, here recede before the original personalities. There is in all independent personality, especially in the personality of genius, something irrational, something new, which does not find its full explanation in what precedes, and which perhaps only slowly finds recognition and adherence. But most of all this holds of the great personalities \opage{14}from whom the new value concepts proceed, whose validity is to be tested through long times in the life experiences of the race. To them applies \textit{the concept of revelation}. This concept is, according to Siebeck, itself a value concept, i.e. a concept whose validity does not rest on the manner in which its content can be demonstrated in reality, but on the practical significance of the content for life. This is seen by comparing different utterances of Siebeck. In one place (p. 192) it says that the religious content that is proclaimed by the great personalities is called revelation on account of its ``supramundane'' character. And as regards the concept of ``the supramundane'', it says in an earlier place (p. 18): ``The concept of the supramundane can first arise on the basis of a world-view that not only takes the world with its existence and content as a given fact, but by virtue of a spontaneous mental activity independent of the given subjects this fact to a \textit{valuation}, the principle of which is drawn from the interior of the mind itself.''

It is really first within this third form of religion that religion and culture come into opposition with each other. Redemption, which in Buddhism is proclaimed in a more negative, in Christianity in a more positive manner, consists in the authority of personality over against the world; it is given with the appropriation of the thought that the world has value only on account of the existence of personal life. Personality has at once emerged from the world and is independent over against the world; it wishes not merely to understand the world and assert itself in it and over against it, but also to work in it for that which, from the standpoint of personal life, stands as the highest. Hence it comes that it falls so naturally to religious faith to conceive the ground of the world itself in analogy with the personal. For faith personality is the first object, whereas it is the very last object of science, because it is the most composite and at the same time most central subject-matter with which knowledge can occupy itself.

On this highest standpoint an opposition arises between faith and knowledge. By virtue of the value of certain phenomena the religious consciousness is inclined to place them outside the natural causal connection, whereas on the other hand \opage{15}science is inclined to declare itself to be the sufficient foundation for view of life and conduct of life. Siebeck maintains against the first assertion that belief in miracles is not a necessary element in religious faith. The clearer the religious consciousness becomes about itself, the more it will --- he holds --- see that belief in miracles does not glorify but precisely limits God, in that it sets nature as a barrier that is to be broken through by God. Against science's assertion of being a sufficient foundation for life he maintains that this assertion itself can only be a matter of faith, since a scientific proof of it is not possible, and since the rejection of the necessity of religion takes place through an act of personality, which is motivated by the presupposition that with science itself the highest good is given. The very need for scientific research rests on this unprovable presupposition\textsuperscript{5}).

There is in all religion a subjective and an objective element. The subjective element, which one might call religiosity, expresses itself in mood, the objective element, religion as opposed to mere religiosity, is thing, work, performance, visible and tangible shape for the subjective mood. The depth of religiosity is the measure of value for the sufficiency and vitality of the objective being of religion; but the subjective mood would never have developed except under the influence of the objective elements, so that religiosity never in time precedes religion. There is a constant relation of reciprocal action between them. But the more a system of cult is organized, the more are new moods and new thoughts excluded from the religious life. It is cult that always makes the greatest resistance against religious renewal. The higher a stage of development religion reaches, the more pronounced individual variations will the religious life present, so that the great human differences can assert themselves not only in the relation between the religions, but also within one and the same religion. ---

This is the most essential content of the first division of Siebeck's work, the descriptive part, which treats of the essence and development of religion. As he now goes on to discuss the truth of religion, he declares that what is here at stake is not the truth of \opage{16}any single religion, but the validity of the religious principle in general. This validity can be grounded only by revelation --- in the sense of the word indicated above --- intervening in human culture and, although it does not itself spring from this, yet affording a necessary supplement to it.

At this point it is not quite easy to understand Siebeck's line of thought. In his descriptive part (which, incidentally, is not merely descriptive) he has sharply distinguished between valuation and causal explanation, so that these need not come into opposition with each other, but denote different ways of viewing existence. But now he sets up the assertion that an ethical action, an action that aims at asserting and carrying through the valuable in the world of reality, becomes possible only if this is in the last instance subject to an ethical power, and he finds in this a ground for the concept of God. Every ethical action rests, he holds, on the presupposition that it is possible to realize the content of the ethical ideas within the connection of nature, and this presupposition is a religious presupposition. Of the ethical power that is here presupposed one can form an idea only by the way of analogy, but such a use of analogy Siebeck regards as justified when it avoids both the dualistic relation between God and world of ``deism'' and the abolition of every difference between them by ``pantheism''; thereby one arrives at a ``theism''. --- Here one misses in Siebeck a closer investigation both of the significance and justification of analogy and of the relation between valuation and causal explanation. As regards the first point, this is the more remarkable in that Siebeck himself earlier in his interesting treatise \textit{Die metaphysischen Systeme in ihrem gemeinsamen} \textit{Verhältnisse zur} \textit{Erfahrung} (in the second volume of Vierteljahrsschrift für wissenschaftliche Philosophie) has demonstrated the significance of analogy for every speculation that goes beyond experience. It remained now only to apply this consideration to the religious ideas. As regards the second point, he does indeed emphasize the difficulties of uniting the teleological and the causal \opage{17}way of viewing in full mutual harmony, but he holds that the religious consciousness does not regard the world as a closed totality, but as part of a higher connection of things, from which ever new creative forces can flow into it. I am not sure that the religious consciousness is here interpreted rightly; it has in any case very often a tendency precisely to regard ``the world'' as a whole for itself, and Siebeck has himself emphasized this in his critique of the idea of miracles. It is rather for the scientific consciousness that the world does not form any closed totality. ``The world'' is not a scientific concept.

In thus acknowledging the justification of conceiving the natural causal connection as subordinate to an ethical principle, Siebeck does not overlook the objections that constantly arise against this way of viewing from the existence of pain and evil in an existence that is supposed to have its ground in such a principle. He does, to be sure, not see that this old objection must strike particularly sharply ``theism'', which conceives that principle in personal form. But he concedes that the difficulty is insurmountable for the religious consciousness, however much this consoles itself with the thought that evil must be possible in a world where free beings were to develop. He has given the chapter in his book in which this question is discussed the heading ``Theodice'' and remarks that theoretically it is God who is to be defended in an attempt at theodicy, but that what matters for the religious consciousness is to keep its hope up in spite of its trials. One might remark upon this that prior to all ``theodicy'' there goes a stage in which it is man who is to be justified before the tribunal of God. Then comes only later a stage in which God is summoned before the tribunal of human thought; then one begins to write theodicies. After that comes the third stage, where neither of these takes place, but where it clearly and distinctly concerns a struggle to keep hope up, and where the ideas one forms of what lies beyond experience expressly stand as symbols for this hope.

But when it thus turns out that the problems that science \opage{18}does not solve are not solved in religion either, indeed appear in it precisely in the sharpest form, then the supplementation of culture, the completion of the world-view, that the religious idea was to bring has not succeeded. The necessity of adopting this supplement cannot be demonstrated.

\afsnit{2. Auguste Sabatier.}{2. Auguste Sabatier}

It is a sign of the times that also from the theological side one begins anew to take an interest in the philosophy of religion. This interest has otherwise been as good as vanished, after the hopes that had been cherished for romantic speculation were shattered. The Protestant-theological professor in Paris who died a few years ago, AUGUSTE SABATIER, has in his work \textit{Esquisse d'une philosophie de la religion d'aprés la psychologie et l'histoire} (5 éd. Paris 1898) set himself the task of treating the religious problem on a psychological and historical basis.

In psychological respect Sabatier lays the weight on this, that what is aimed at in religion is not a solution of ``the riddle of the world'' but of ``the riddle of life''. It is the enigmatic and contradictory in our own being that gives rise to the religious need, and world conditions interest from the religious standpoint only so far as they are assumed to have influence on the solution of the riddles of the inner life. The religious need is a kind of higher need of self-preservation, which is awakened by the disharmonies in our inner life, especially by the opposition between dependence and freedom, or between passivity and activity. The stream of our life, which springs from the innermost in us, is broken and split like a powerless wave against the skerries and rocks of external things. This is the ground of all pain, and then especially of the highest pain of all, that which is conditioned by the disharmony between our ideal and what we can carry through in reality. A harmony between these oppositions is attained, according to Sabatier, in prayer, the proper religious act. Here there appear at once the feeling of dependence toward the mysterious power we feel ourselves related to, and the feeling of freedom by virtue of the force that streams to us from this power. The essential in prayer is not the words in which it \opage{19}is clothed, or the ideas to which it is attached, but the inner connection of self-surrender and self-activity, which here mutually support each other and lead over into each other. The necessity of religion rests on its springing from our innermost conditions of life. Therefore it will show itself even where there is antipathy against every external religion. There is an inner altar that can be preserved and kept in honor even if one overthrows all external altars. There is only one true atheism: frivolity and blasé indifference.

When the need is strong and inward, the satisfaction is already given; it does not come from outside. For Sabatier the whole mystery of religion lies in Pascal's words: ``You would not seek me if you had not found me''. The answer, revelation, is present in the very possibility of prayer, in the inner force and striving from which prayer springs\textsuperscript{6}). At the same moment man seeks God, it is God who seeks man. Revelation is not a communication once for all of unchangeable doctrines, which it then only matters to hold fast, but consists in the manner in which the consciousness of God arises and develops in the single individual and in humanity.

Sabatier designates this concept of revelation as the psychological or the critical, in opposition partly to the mythological, partly to the dogmatic concept of revelation. It has its great historical representatives in the Israelite prophets and in Jesus, but it finds application in every religious personality. Religious inspiration is as mysterious as the poetic, but no more miraculous than this. It is an inner fertilizing force, which works in harmony with all mental forces and laws. Prophetic inspiration is only a higher degree of what stirs in every strongly moved religious state of feeling. The standard for the value of revelation lies in the enrichment, liberation and elevation of the personal life that it effects.

The ideas and actions in which religion finds expression vary with the times. It is the great fault of orthodoxy, even of the Reformers, that they have not distinguished between the religious experience and its expression. The psychological phenomena and the theological explanations have been confounded. In \opage{20}every dogma a distinction must be made between the religious experience that lies at its ground and the more or less imperfect expression that reflection has been able to give it. This last element is the most variable element in religion. There may be three kinds of dogmatic variation. Dogmas may decline and disappear for want of use (par désuétude), in that religious experience no longer affords the content they are to express. Thus it goes with the assumption of demonic possession and of a personal devil. Dogmas may next acquire another content than the original, in that they are interpreted according to new experiences. Such a displacement goes on uninterruptedly. By inspiration, incarnation and Trinity later generations do not understand the same as earlier generations. New wine is constantly poured into the old wineskins. Finally new dogmas may arise as expressions of new experiences. Thus there arose within Protestantism justification by faith and the universal priesthood.

To proclaim the immutability of the dogmas of the Church would be to proclaim the immutability of the philosophy that was employed to shape them. Dogmatics must keep to the religious experiences of which it is the exposition. And when certain concepts are employed in order to be able to make this exposition, it must also acknowledge the results to which the theory of knowledge comes regarding the significance and validity of our concepts.

Religious knowledge is symbolic knowledge. To express the feelings of trust, fear and love it uses images that are drawn from human relations. The power of the symbol rests on the close connection between feeling and image. Religious feeling does not at first distinguish between content and form; it is not conscious that it expresses itself symbolically. It is critical philosophy that impresses the unavoidability of symbolism, which both orthodoxy and rationalism each in its own way misjudge. --- The symbols may become common to large circles, and there are symbols that perhaps will never disappear from the religious life.

Not only the theory of knowledge does Sabatier require to be respected by dogmatics, but also the general method of the empirical sciences. He distinguishes, as we have heard, between the riddle of the world and the riddle \opage{21}of life. The riddle of the world it is the business of science to discuss. The world must be taken as science shows us that it is, and the riddle of life is then posed by the question of how science's image of the world relates to the innermost need of the spiritual life --- whether it contains means or hindrances for the satisfaction of this need.

The independence of empirical science must benefit psychology too, and here we stand (as I have already pointed out in my Philosophy of Religion p. 173) at the most difficult point in Sabatier's philosophy of religion. He has traced all phenomena of the religious life back to religious experience, as he describes this more closely, and he has maintained that all religious ideas are symbols that express this experience. Yet he nevertheless assumes that God is the cause of religious experience, and here there would then be a case in which a phenomenon of experience received a theological explanation, the possibility of which Sabatier himself has expressly rejected, since God is not himself a link in the series of phenomena and therefore cannot be used for scientific explanation. Moreover, the very idea of God, like all other religious ideas, must surely have a symbolic character\textsuperscript{7}).

Sabatier's theological point of departure shows itself, however, not only in his attempt to explain religious experience but also in his very description of it. He brings forward only a single type of such experience and does not lay stress on the variations, although in principle he acknowledges their possibility and their legitimacy. As with Siebeck, we feel here too the lack of a broader psychological material. Such a material would perhaps have led to a less narrow determination of the nature of religion. Yet these objections do not prevent one from parting with the French theologian's philosophical work with the memory of a noble, refined and enthusiastic man.

\afsnit{3. Rudolph Eucken.}{3. Rudolph Eucken}

After a series of preliminary works (especially \textit{Die Einheit des Geisteslebens} 1888, and \textit{Der Kampf um einen geistigen Lebensinhalt} 1896) this brilliant thinker, who works in \opage{22}Jena, the old hearth of idealism, has expressed himself on the religious problem in his book \textit{Der} Wahrheitsgehalt der Religion (1901).

It is not a system of philosophy of religion that Eucken has here wished to give. Such a system is, in his view, for the present impossible on account of the unclear and uncertain conditions in the religious domain; ``In the spiritual anarchy of our age there is no fixed and recognized point of attachment; every more deeply penetrating discussion must go back to the first foundation and build up anew upon it. Thus I too had, through a general investigation of human existence, to work my way step by step to the place where the problem of religion breaks forth, so as then immediately to appear as the center of all the search for the soul and significance of our existence''. The author declares himself a seeker, and he addresses himself to seekers. He finds in the age a strong longing for religion combined with a distinct consciousness of the inadequacy of religion's present form. The new type of life that the Renaissance asserted, and the new labor of thought thereby initiated, led away from Christianity. But the grandiose culture in which this type of life and this labor of thought took shape has brought with it inner entanglements, and faith in its sufficiency has been shaken. Thereby the thought emerges whether there might not be something in religion that cannot and must not be lost. The lack is felt of a great goal, which not only could cast light over the great circle of human life but also raise the single individual above his often petty circumstances. And at the same time the old riddles of life, which one regarded as solved or at least as displaced, press forward anew. Above all a sharp contradiction is felt between man's spiritual predisposition and his actual condition. All this leads to a struggle between religion and culture.

But neither religion alone nor culture alone commands the means of deciding the struggle. The decision must come about by distinguishing, both within religion and within culture, between the eternal and the transitory, and between the indispensable and the intolerable. A revision must be undertaken of the very fundamental conditions of our spiritual life, and to \opage{23}undertake such a revision is the business of philosophy. Not in the sense that philosophy can construct a religion. For religion is something factual, an expression of a real existence, which can be pointed out when it is there but not artificially produced. But since what is at issue here is the pointing out of a spiritual fact, the pointing out is not a quite simple matter; spiritual labor is needed to find religion's peculiar place in the whole connection of life. Philosophy cannot get further here than to point out possibilities, which only through great personalities and often under great upheavals in human conditions can be converted into actuality. But if it can be demonstrated that religion is bound up with the innermost ground of our existence, then every attack upon it and every criticism of it can serve only to make its innermost, eternal core more and more manifest.

The problem is posed by the position of the spiritual life in existence. Its persistence and its development presuppose that behind and beyond the spiritual processes that experience shows us there lies an eternal and coherent spiritual world (eine geistige Substanz, ein seelischer Bestand), which contains the \textit{possibilities} for new formations and the \textit{yield} of what has already developed. The life of the mind must be more than a mere phenomenon, hence more than psychology can point out. It must have an eternal principle in itself. Not merely a psychological but a metapsychical view must be adopted. For such a view essence and value are one, in that its highest object is the valuable core of existence itself. Only through such a view does a central content of life become possible.

The ancient culture overvalued form; the culture founded by the Renaissance overvalued force. The true culture, which has a religious character, sets \textit{formation of essence} (die Wesensbildung) in contrast both to formation of form (die Formbildung) and to formation of force (die Kraftbildung). Both form and force are to serve a valuable content. Formation of form alone leads to a congealing in plastic works of art, in which the eternal is supposed to be expressed for all times. Formation of force leads to restless movement, so that all being is dissolved. Formation of essence, \opage{24}on the other hand, maintains a constant interaction of eternity and time, in that the eternal possibilities are converted into actuality through labor in time, and in that, on the other hand, the yield of this labor is preserved in the form of eternity. Preservation of the given and production of new content must go hand in hand, in order that we may press forward to true reality.

The line of thought that maintains the possibility and necessity of ``formation of essence'' Eucken calls the \textit{noological} (from noos, reason, mind, which among the Neoplatonists is placed above psyche, soul), and he sets it in contrast to the psychological view, which merely concerns itself with the conscious life given in experience. It is, in his conviction, nothing new that he thereby proclaims, but an old truth, which KANT in particular has impressed by his definite distinction between the logical, the ethical and the esthetic mode of consideration on the one hand, the empirical-psychological on the other. Here validity or value stands in contrast to reality --- but in such a way that validity is in the end maintained as the highest reality. --- The noological line of thought differs from the speculative or metaphysical line of thought of an earlier age, which believed it could attain an understanding of the phenomena by theoretical abstractions. The noological method builds on experience, on fact. To explain noologically means to place a spiritual manifestation of life in the whole of the spiritual life, point out its position and its task within this whole, and thereby illuminate and consolidate it. This explanation differs from the psychological, which consists merely in pointing out the processes through which man acquires and appropriates a spiritual content. Both methods must be applied, in order that the opposition between idealism and realism may be overcome. Idealism has maintained the independence and intrinsic value of the content of mind; the justification of realism rests on its emphasis on the natural conditions. The realistic mode of consideration comes into its own where it succeeds, within the spiritual whole that the noological conception maintains, in finding the special causes and laws; but it is wrong in thinking that causal explanation can furnish the last standard of truth. Idealism has rightly demanded a \opage{25}valuable content for existence, but has disregarded the real conditions.

According to Eucken, the noological view now leads precisely to the fact that religion maintains: the presence of an absolute life of the mind, exalted above and yet at work in the phenomena of life appearing in experience. There is no religion without the living presence of a higher world in our empirical world --- but in such a way that these two worlds stand in an irrational relation to each other. Religion can exist without belief in God; genuine Buddhism is an example of that. But without an outlook beyond the world of experience religion is an empty word. The motive for religion lies in the opposition between essence and form of existence: in mere experience the life of the mind stands as an appendage to the material, as a secondary matter; but this conflicts with the unconditional demands that the life of the mind makes, and the exalted values that it maintains.

From the fundamental claim of religion, reality looks different from before. A view of the world proceeding from religion will direct its eye toward what abides (in contrast to what changes), toward the free (in contrast to the merely natural) and toward the rational (in contrast to the irrational). It will take particular note of the points where a new sphere begins, such as the transition from the inorganic to the organic, from the life of nature to the mental life, from the mental life to the life of the mind, and it will conceive this graded series as an unfolding of the innermost in existence. But such a religious conception of the world cannot be proved. It has the character of a conviction, not of a science. It is the business of philosophy to determine the relation and the boundary between the scientific and the religious view of the world.

The motive for religion lies in a spiritual need. It is not a question of man's empirical felicity and its persistence, but what is decisive is the wholeness and unity of the life of the mind, the foundation of all special and individual manifestations of life. It is its persistence that is at stake, in that mere experience does not confirm it. It then depends on whether there can be awakened a strong need to hold fast a conviction on this point. Empirically viewed there are grounds enough for doubt. Experience \opage{26}by no means shows us the world of the mind as an independent world, still less as the main thing in existence. And in the course of the development of culture personalities are sacrificed to society and to progress. Within the life of the mind itself disharmonious tendencies stir; part stands against part, and part against whole; spiritual tendencies combat one another, and evil proves to be a real power.

In face of the doubts that thus arise, Eucken stresses, first, that the strong feeling of disharmony is itself a consequence of the great demands made upon life. The bitterness of suffering and the harshness of resistance bear witness to the depth of our existence and to the presence of higher forces in our being. With the oppositions pain too would fall away. It was never in times of unhappy experiences and heavy entanglements that mankind doubted its ideal task, but on the contrary this doubt thrived in such times as sluggish comfort and imagined satiety prevailed.

Second, it is Eucken's conviction that that doubt must be combated by a more powerful and more concentrated form of religion than the indeterminate and universal religion that is described in what precedes. He points to the significance of the great personalities in the history of religions. In them the divine stirred in a deeper and more powerful manner than in others, and they were able to express life's manifold experiences in clear and glowing images. A new life broke through here in a decisive manner and intervened decisively in the historical development of the race. By bringing forward their significance we make the transition from \textit{the universal religion} to \textit{the characteristic religion.}

``Characteristic religion'' denotes in Eucken the historical or positive religions, in so far as they --- each in its own peculiar manner --- give a new total image of life, clear in its fundamental features, and set this image in contrast to the world of experience. It differs from ``universal religion'' by the determinate shaping without which such an opposition to the world of experience cannot be maintained. It springs from great personalities and rises above all popular culture as an immediate utterance of the innermost essence of the life of the mind. Without \opage{27}such characteristic religion independent inwardness is not possible. It is here not a question of a mere embellishment of life, but of the breakthrough of a new life, a new reality. Yet the breakthroughs to which the great positive religions bear witness are only high points and turning points within a great process. The divine, which here breaks through, did not first come into being at these particular places.

The higher life that thus breaks forth our concepts are not able to express. Imagination must come to our aid with its symbols. And here lies religion's vulnerable point. Imagination is always determined by the given stage of culture and its traditions, and the symbols will therefore acquire a historical character. A time will one day come when that which perhaps for millennia was an expression of the highest truth, and which for human consciousness had entered into an intimate connection with the religious content, is shaken in its significance and must yield its place to new forms --- often only after severe tremors and hard struggles, during which doubt, from being directed against the symbols, may pass over to attack the innermost core of religion.

The universal religion is constantly necessary as a corrective of or introduction to the characteristic religions, which, left to themselves, are easily narrowed and become mere means of consolation. The significance of the characteristic religions rests not on the fact that any of them should have found a conclusive expression of the highest truth, but on their vivifying and elevating effect amid the unrest and vicissitudes of the time. It must ever anew be tested whether what is historically handed down really performs such a work. But the Church has a tendency to place itself slavishly before a model given once for all and to conceive all truth as consisting in imitation and repetition. It isolates the highest and thereby gives it a passive position in relation to real life. Thereby struggle always becomes necessary in the religious domain. It is the significance of the great religious personalities that they compel us to choose. ---

What especially determines Eucken's place in modern philosophy of religion is his decided conviction that the \opage{28}values can be maintained only when, behind and beyond the world of experience, a higher reality is assumed, in which the life of the mind that in experience appears scattered and in changing forms persists in the form of unity and eternity. In an often profound and atmospheric manner he maintains the necessity of this opposition, --- now so that he recalls the old mysticism, now so that he recalls romantic idealism, and now so that it is the exalted poetry of the personal struggle of life that finds expression in his words. Yet he stresses very strongly that philosophy can here only point out possibilities, and that it is through prophetic personalities that such possibilities gain actuality and definite form. When he --- in order to avoid subjectivism --- maintains the necessity of a metaphysics (or a ``metapsychics''), it is yet properly by virtue of the \textit{value} of the true, the beautiful and the good that this metaphysics is to be necessary. What he calls the ``noological'' method differs from the psychological method in that it will show that that value can be maintained only by the assumption of a connection of the whole that bears the valuable manifestations of life appearing in our experience scattered and fleetingly. In particular the noological mode of consideration lays stress on the great breakthroughs, the significant beginnings and the qualitative differences within the life of the mind as we know it in experience. He regards the attempts of empirical science to show continuity in the world of phenomena as hopeless, but interprets the breaks in empirical continuity as single outcroppings of a great transcendent continuity, a realm of possibilities, ends and values, which only by glimpses reveals itself in the world of finitude. In this he is akin to romantic idealism. But noology differs from the speculative method in that the assumption of such a connection of the whole is expressly grounded only on experience of the valuable. It seems then clear that the point of departure is empirical and personal. He gets his impetus to pass over from experience to ``metapsychics'' through the tension that is produced by a great dilemma. Either --- he maintains --- one must regard the life of the mind as a fleeting phenomenon, an appendage to the material \opage{29}nature, --- or one must assume that it in itself constitutes a great totality and is one with the innermost core of existence. This dilemma comes forth in him both where he develops the idea of ``the universal religion'', thus would make the transition from the empirical standpoint to a religious standpoint of a more indeterminate kind, and where he goes over to ``the characteristic religion'', thus would show the necessity of a historically formed religion. Characteristic are especially the following utterances: ``In this decision it is not this or that but the life of the mind as a whole that is at stake. In man's experience the life of the mind appears in an untenable intermediate position; from this man must either turn back to mere nature and confess that all peculiar striving was only a heavy dream, or boldly go forward and secure for the new striving a new world. Only an unclear way of thinking can seek a middle way between these ways, and only a feeble disposition can calmly endure the dilemma. In any case it stands fast that without religion there is no truth in the life of the mind and no inner greatness in man.'' (Der Wahrheitsgehalt der Religion. p. 238.) ``Is human life a mere appendage to nature, or is it the beginning of a new world? On the answer to this question depends all shaping of our relations, the whole direction of our action. Religion would raise human existence up to a height raised above the world and thereby save our life from the nothingness into which it otherwise surely falls. If this its attempt is only a flight of Icarus, then all hope falls away, then the noblest and best stands as empty imagination, then the whole ends in unreason.'' (ibid. p. 319). ---

What is to drive us to choose here is nevertheless always in the end personal experience of life, in particular the degree to which one is gripped by the fate of values in reality. There seems to be no ground for introducing a new term; ``noology'' is neither a necessary nor a happy enrichment of the already sufficiently supplied philosophical terminology. From a private utterance of Eucken's, which I hope to be able to cite without committing an indiscretion, it is seen clearly that it is the personal experience of life's \opage{30}oppositions, in particular of the great opposition between ideal and reality that has determined his own choice and led him to ``noology''. In a letter he wrote to me after reading my Philosophy of Religion, he states, after having dwelt on the predominantly psychological character of my attempt to treat the problem: ``Mich treibt es, bei allem Gegensatz zur alten ontologischen Metaphysik, immer wieder zur Metaphysik und so auch zu einer metaphysischen Begründung der Religion; ich fürchte ohne das dem Subjectivismus zu verfallen, und dagegen sträube ich mich namentlich deswegen, weil ich stärker unter dem Eindruck der grossen, in der Erfahrung unlösbaren Konflikte unseres Daseins stehe, kurz gesagt, dualistischer denke und denken muss.'' In this utterance Eucken points directly to individual experience as determined by the peculiarity of the single personality. It is this that determines his standpoint. But the more clearly this foundation appears, the more natural it becomes to assign to the philosophy of religion the task of investigating the psychological soil from which the different religious standpoints grow up, and to attach to this an examination in the theory of knowledge and an ethical examination of the thoughts and endeavors in which these standpoints find outlet. A quite peculiar method the philosophy of religion then does not need, and the so-called ``noology'' is indeed itself a ``middle way'' between speculation and practical faith. As speculation it falls to the treatment of the theory of knowledge, as practical faith to that of psychology and ethics. The conception of the task and procedure of the philosophy of religion that follows from this seems to be most in the spirit of critical philosophy, in which Eucken --- like Siebeck and Sabatier --- would work. The ``metaphysical'' grounding of religion that Eucken demands of philosophy proves, according to his own utterances, in the end to rest on the need for a decision of a great practical dilemma.

I find in Eucken's standpoint --- in the ``noology'' that leads him to his religious metaphysics --- an example of how personal experience of life can lead to faith in the persistence of value. For me every faith, even though it appears in the most philosophical forms, is always only \textit{object} of philosophy, not \opage{31}itself \textit{product} of philosophy. In any case the philosophy that produces it is itself a product of personal experience of life. Philosophy as science has, after all, no other methods than the logically inferring, the psychologically explaining and the ethically evaluating.

But from my standpoint such works as Eucken's, and the same holds of Sabatier's, acquire an exceedingly great interest through the clear manner in which the decisive motives for a religious view of life come forth. Two men of ideal mind, gripped by life and equipped with energy of thought, have here expressed themselves on one of the greatest human problems.

\kapitel{II.}{Investigations in the Psychology of Religion.}{II. Investigations in the Psychology of Religion}\opage{32}

A thoroughgoing attempt to go back to immediate experience in treating the religious problem must in particular direct attention to the point where consciousness passes from immediate experience to forming a judgment, thus to the first setting up or involuntary coming into being of a religious proposition. What goes on here cannot be essentially different from what always goes on in the transition from observation to judging. Religious consciousness has neither a psychology nor a logic that differs from the ordinary psychology and logic, and religious judgments must come into being as all other judgments come into being. What is peculiar to religious judgments is that the experience on which they ultimately rest is an experience of feeling, an experience of the value of certain experiences as furthering or hindering, and that the experiences that are here in question concern the relation between ideal and reality, between the world of values and that of reality. Religious values are therefore values in the second power or secondary values: they presuppose that man knows something that gives life value (in whatever this then consists), and that he then experiences that there are powers in the world that further or hinder this valuable thing. At the high points of religious feeling the difference between primary and secondary value naturally falls away; there is known no Why and no predicates are applied. In most pronounced form this occurs in mysticism. But where a judging can take place, thus where there can be talk of statements of religious consciousness, there --- \opage{33}I maintain --- a closer analysis will show that what something is said about is the relation between value and reality, as the individual believes he has immediately experienced this. Religious judgments concern the fate of values in the world of reality. --- Here it must, however, well be remarked that religious judgments are not always, perhaps even very seldom, formed at first hand.

In every formation of judgment one must (as I have more closely developed in my treatise \textit{Det psykologiske Grundlag for logiske Domme}) distinguish between the initial idea and the concluding idea. In the linguistic formulation of judgments these two ideas appear as logical subject and logical predicate. At the beginning of the process of judgment the initial idea has a certain indeterminacy, which is sought to be removed by the discovery of a predicate, a closer determination, which takes place through analysis of the initial idea and of the connection in which it occurs. The richer the content of the initial idea is, the more difficult it will be to find a definitive concluding idea, a concluding and adequate predicate. In this respect religious consciousness, as it appears in history, presents the greatest differences. The narrower both the world of values and that of reality is, and the more elementary the relation between them appears, the easier it is for man to rest content with a predicate and thereby conclude religious thinking. Here are all possible degrees from the limited and childish predicates of fetishism and animism to the sublime limiting concepts of the highest mysticism. And the process by which the concluding idea is here linked to the initial idea, is now involuntary association of ideas, now imitation and habit, now speculative brooding, now anxious search and questioning. The content of the initial idea in religious experience is the harmonious or disharmonious relation between value and reality, and the concluding idea is to state what it is that comes to light in this relation. If such a concluding idea is won or found, the transition has been made from experience to faith. The concept of God in its various forms is the basic predicate in which religious consciousness seeks to gain expression for and explanation of its deepest and most essential \opage{34}experiences. The word God is to mean ``him to whom sacrifice is made'', thus him to whom man feels himself in debt, to whom he refers his experiences of the relation between value and reality. But who is that? Qvid est quod amo, quum te amo? asks Augustine. Here lies the possibility of ever new brooding, especially since the initial idea is not constant, inasmuch as the relation between value and reality is constantly changing.

There is in Goethe's Faust a monologue that suits well as an example of what is here meant by religious experience. In the scene ``Wald und Höhle'' Faust makes a reckoning. Nature's and thought's great treasures are opened to him; he can immerse himself in the inner as in the outer world in understanding and mood; but at the same time there has been given him a demonic companion whom he cannot do without, and whose coldness and impudence transform those high gifts into nothing --- at the same time as he awakens the flaming fire of sensual desire. A concluding idea does not come about here; if one wishes to seek one in the poem, it must be found partly in the Prologue to the first part, partly in the end of the second part. --- The great opposition that Goethe's Faust here accounts for lies --- only in still sharper form --- at the ground of Schopenhauer's experience of life, and Schopenhauer's philosophy was in its innermost ground an attempt to find a concluding idea as expression for this experience. ---

Immediate religious experience is an ideal, since every human being we know begins under the influence of a tradition and at once interprets his experiences by the help of tradition, or perhaps more exactly, at once sees them in the forms and the illumination that tradition has arranged. Already the growing into a language brings about that there can here be no talk of absolutely original experience. Every religious experience is more or less confessional. Both the values, reality and the relation between them man apprehends from the first under the influence of ideas that he has not himself formed on his own, still less himself derived from his experiences. It will then be of great interest to investigate how the traditional ideas are gradually altered by the individual's own advancing experiences, and next to make a comparison between religious \opage{35}standpoints and processes of development that are determined by essentially different points of departure.

There now exists on all the topics touched here --- both on religious judgments as value judgments, on the forms and modifications of a confessional religiosity and on comparative psychology of religion --- a series of works, which for the most part have appeared after the completion of my Philosophy of Religion, and which it will be of interest to have characterized.

\afsnit{1. On Religious Judgments. {\normalfont\textit{(Max Reischle).}}}{1. On Religious Judgments (Max Reischle)}

The attempts at philosophy of religion with which we occupied ourselves in the first chapter were more or less consistently built on critical philosophy and its distinction between purely theoretical judgments and value judgments. Within theology this distinction has in recent years played no small role, especially with ALBRECHT RITSCHL and his disciples. Auguste Sabatier stands close to this tendency, and he also speaks of value judgments (jugements d'estimation ou de dignité). The weight that is laid on the difference between the two kinds of judgments is characteristic of the endeavors of the most recent time to lead religion back to immediate experience and to determine the proper character of this experience. It is through the investigation of value judgments that in the most recent time there has come about a contact between theology and philosophy, which otherwise in the last half-century --- especially on account of theology's more and more pronounced confessional character --- have stood far apart from each other. --- I shall here as an example draw attention to a treatise by MAX REISCHLE: \textit{Werturteile und Glaubensurteile}. (Halle 1900).

A value judgment presupposes, according to Reischle, a feeling of value that differs from the mere feeling of pleasure. The condition for the arising of a feeling of value is that the satisfaction concerns the whole self (das Gesammt-Ich), not a single part or state of the self. Value I attribute to a thing of which I have seen by way of reflection that it will satisfy my self as a whole. Yet Reischle concedes that the boundary between ``the whole self'' and its single parts or states is fluid. What he wishes \opage{36}to stress is especially that before the value judgment there goes a composite process, by which the feeling of value attains a clarity and certainty that the mere feeling of pleasure lacks.

Value judgments can, according to Reischle, be of two kinds: natural and ideal (or normative).

\textit{Natural value judgments} are grounded in man's nature and conditions of life, which determine the satisfaction that is possible for him. Here many differences make themselves felt between the different individuals, but since interests are nevertheless to a certain degree solidary, there can be collective value judgments, and since conditions of life are to a certain degree common, there can be general value judgments. Where an external standard (e.g. money) is possible, the value judgment can acquire an objective character, so that perhaps even its subjective origin is forgotten.

\textit{Ideal or normative value judgments} presuppose an idea or a norm that we recognize and expect to be recognized by others. Here belong esthetic, ethical, intellectual and religious value judgments. These are determined not merely by the manner in which we are passively determined from outside, but essentially by agreement with a law that holds for the activity of our own mind. In ethical value judgments e.g. we follow a norm or an idea that is able to lead our will to harmony and freedom. In religious value judgments the underlying idea is the significance it has to submit oneself to a highest power in order by unity with it to reach the highest freedom in face of the world. The ideal value judgments make claim to general validity, because no one can deny the ideas that condition them without giving up his personality, the freedom and inner connection of his being. ---

The difference between natural and ideal value judgments is --- what Reischle has not noticed --- evidently just as fluid as he himself concedes the difference between feeling of pleasure and feeling of value to be. That ``self as a whole'' which the feeling of value is to presuppose, namely contains a norm that already holds for the special and changing elements and states; their value is determined by the relation to it. And besides, on the other hand, ``the self as a whole'' must itself demonstrate its value in face of the single elements and states whose subordination is demanded; \opage{37}it is by no means self-evident that the whole is in the right against the single element or the single state. Finally, according to Reischle himself the value of the esthetic, ethical, intellectual and religious functions of mind rests on their procuring our being harmony and freedom; --- but that means nothing other than that they make a ``self as a whole'' possible, and the ideal value judgments become, seen from this side, not different from the natural. It is therefore by no means easy to set up measures of value in the spiritual domain, where there is the most intimate interaction and mutual determination of the elements among themselves and between elements and whole. Within the individual consciousness, as within society, the different measures of value struggle. The value of every idea or norm must itself be demonstrated --- and therefore the problem is posed ever anew. Of all value judgments precisely, as I have shown elsewhere, those that lead to the setting up and recognition of a norm present the greatest difficulties for theory\textsuperscript{8}). The transition from the ``natural'' to the ``normative'' (ideal) includes great problems, which a mere classification of the actually given value judgments easily covers. Reischle has a clear eye for the analogous problem that arises when the different religions, each with its basic values, stand over against one another; but he has not seen that the struggle can already begin in the interior of the single individual, and that a guiding idea or norm is always the result of such a more or less violent struggle. ---

In the last part of his investigation Reischle discusses how far \textit{articles of faith} can be called value judgments. He concedes that authority and tradition play so great a role both in the Protestant and in the Catholic Church that the articles of faith that the single individuals profess cannot be called value judgments in the sense that these single individuals themselves have derived them from their experiences of value. On the other hand Reischle considers it undoubted that the articles of faith must be called value judgments when it is asked on what their claim to validity rests in principle. They are value judgments of ideal or normative kind; their significance consists in their stating the conditions for blessedness, i.e. for the personality's complete harmony and freedom. It would be wrong to call them postulates, as if it were merely our \opage{38}own need and yearning that formulated in them what could produce satisfaction. There lie at the ground of the religious judgments definite facts, for the Christian thus the life and death of Jesus, and the religious judgments express trust in the power and significance of such facts\textsuperscript{9}). The Christian believes in Christ, because only in adherence to him can he find the true life. Of all historical religions it holds that certain given phenomena --- persons and events --- are interpreted as divine revelations, and the religious judgments express trust in these revelations. Striking, especially beneficent or destructive phenomena in nature, in the conscious or the unconscious mental life, in the people or in history awaken the involuntary idea that higher powers are here at work for the weal or woe of men. Often the judgment is dictated by fear, most often however by trust. And the higher a religion stands, the more it is ethical ideas that dictate the judgment. The very fact that something is called ``revelation'', is thus according to Reischle a value judgment; it depends on whether the personalities or events in question stand as forms of the entry of great values into history, and whether they can become object of trustful adherence and devotion. A purely objective decision of the strife between different revelations is therefore impossible. This is a necessary consequence of empiricism in the philosophy of religion, and Reischle also draws this consequence, which he denies only by putting the word ``revelation'' in quotation marks where it is not a question of his own religion.

\afsnit{2. The Pathology of Religious Feeling. {\normalfont\textit{(Murisier).}}}{2. The Pathology of Religious Feeling (Murisier)}

To be able to point out the elements in religious feeling it may be expedient to study its extreme forms, in which it assumes a morbid character. The morbid consists in elements that otherwise occur in harmonious connection with other elements separating out and appearing with overwhelming strength, so that they dominate the whole course of life. In his book \textit{Les maladies du sentiment} religieux (Paris 1901) E. MURISIER (Professor in Neufchâtel) has set himself the task of \opage{39}investigating the morbid forms of religious feeling, which he thinks he can trace back to two, ecstasy and fanaticism, the former the extreme expression of religious feeling in its individual form, the latter its extreme expression in social form. These two forms stand in opposition to each other: the one leads to absolute isolation, the other to outward-going activity and to connection with others. But there lies, as Murisier shows, one and the same need at the ground of both types. The mystic sees in isolation and in the removal of all multiplicity from consciousness the only means to preserve unity and security, to hold fast life's total direction toward a single goal. His mind lacks energy to master the multiplicity of life in him and outside him, and he therefore collects himself on a single point. The fanatic seeks, on the contrary, to attain unity and security in life's course and direction by a practical acting that leads as many as possible to share his faith. He seeks to bring about a social unity within which no disharmonious and disturbing tendencies are felt. By social uniformity the unity and identity of one's own self is supported. The fanatic praises ecclesiastical unity with just as great zeal as the mystic the unity of thought. And just as eagerly as the mystic keeps away all impressions that could disturb his absorption in the One, which is the highest for him, just as eagerly the fanatic suppresses every tendency that could break the social unity. His persecuting corresponds to the mystic's asceticism, and his proselytism to the mystic's concentration. For many people adherence to a religious community, especially a narrow and exclusive one, is the only way along which they can find equilibrium and peace of mind.

Through the study of these opposite forms of religiosity, supported on a historical, biographical and medical material, Murisier finds in the character marks that are common to them a contribution to the characterization of religion. Religion is due --- whether it appears in individual or in social form --- to a need for unity that is closely connected with the need for self-preservation. This need for unity appears in religion particularly as a need for guidance (besoin de direction), and in so far Schleiermacher's description of religion as feeling of dependence was a \opage{40}happy stroke. In morbid form the direction toward unity, the concentrated dependence can be attained only through exclusion of heterogeneous impressions and directions; in healthy form it is attained in a more harmonious manner through ``edification'' and organization.

Religion in its social form goes before individual religiosity; historians of religion therefore concern themselves mostly with the former and are inclined to regard the latter as a sign of dissolution. The religion of primitive peoples consists mostly of ideas of faith and ceremonies that are common to the whole society. The self of primitive man has not separated itself out from family, clan and people, just as the child's consciousness still draws its first content from the nearest surroundings. The very ecstasy of the savages' sorcerers reminds more of the fanatics' state of mind than of the mystics' ecstasy; a certain development is required, before the properly individual form can come into being. But the sociology of religion can at no point be quite separated from the psychology of religion. The very need to produce social unity or join one has a psychological cause, which must be studied in the single individuals. Even where religion most nearly fulfills a social demand, it is yet essentially borne by an inner yearning in the minds of the single individuals. ---

Much can certainly be objected to Murisier's conception of mysticism as a purely individualistic phenomenon. Partly the mystic begins as a rule within a confessional religion, and partly his own self belongs to that from which he wishes to be freed; he wishes to get out of the narrow conditions, beyond all limits, and mysticism is often an altruism driven to the extreme. Mysticism is a phenomenon of expansion no less than a phenomenon of concentration, or rather, there are manifold nuances within the mystical direction\textsuperscript{10}). --- There are also more forms of morbid religious feeling than the two main forms to which Murisier keeps. There are thus states of anxiety and despair, where neither expansion nor concentration is possible, and there is in contrast to this the exalted self-feeling, where honor and power of being linked with the Highest are enjoyed in an egoistic manner. This opposition between potentiated feeling of impotence and potentiated feeling of power is perhaps no less \opage{41}significant for the pathology of religion than the forms described by Murisier.

What does not come forward sufficiently in Murisier is the connection of religious feeling with experiences of value. One could have a guiding thread for finding the morbid forms of religion by proceeding from the presuppositions to which the study of religion will again and again lead back: experiences of values, of realities and of the relation between values and reality. But the most scientific standpoint would nevertheless be that which regards the pathological forms of religion as special cases of the ordinary forms of mental disease. Taking up this standpoint would also lead to seeing the limitation within which one must move when one from the phenomena of the sick conscious life would draw conclusions concerning the nature of the healthy conscious life. --- These objections nevertheless do not prevent one from acknowledging that in Murisier's little book there is found much good and fine psychology, and that it --- within the limitation the author has chosen --- illuminates the religious life from a characteristic side.

\afsnit{3. Contributions to the Psychology of Confessional Religiosity.}{3. Contributions to the Psychology of Confessional Religiosity}

From the religious side it is maintained more and more that the foundation of religion is experience, and every attempt at a psychological investigation of religion leads to the same result. But then it becomes of decisive importance to investigate the kind and character of the experience that can be in question here. Otherwise one cannot know how much this foundation is able to bear, and dogmatism once again makes its entry. But here a difficulty now arises, which comes from the fact that every state of feeling has a tendency quickly to enter into connection with definite ideas and to lead to definite actions. And the ideas and actions that always lie nearest at hand are those prevailing in the surrounding society. The individual lives himself into them, perhaps even before he has begun to make his own life experiences. He inherits the valuations of the race before he can begin to value life on his own. And when he \opage{42}begins to make his own experiences, he will in any case at once settle them into the received forms. To a large extent, if not for the most part, his experiences are called forth precisely by the impression that received ideas and actions make upon him, especially through the captivation of those around him by them. Only with great difficulty and under special conditions will that counter-test become possible which is needed for a real proof from experience, --- namely the test of whether the experience and the valuation of life to which the individual himself, according to his nature and his circumstances, can attain, can now really be expressed and explained solely and exclusively through the received forms, so that content and form wholly correspond to each other. Possibly the same content could be expressed through other forms, or other ideas and actions could even make a richer and more valuable life experience possible.

Religious experience therefore becomes, historically viewed, always more or less confessional. It becomes predominantly a repetition of what others have previously experienced and lived through; it moves in ready-made expressions and formulas, and that which at the very outset, namely when the expression was found and the formula was created, was effect and result, operates in confessional experience as cause, often the chief cause. Yet within confessional experience there are all possible degrees between the passive living into given forms and the independent appropriation, which at its height can go on to find new paths and create new forms.

The psychology of religion must keep this fact in view when it gathers its material of experience. A rich material we have here --- besides what one's own knowledge of human nature yields --- in the utterances of religious personalities, especially in autobiographies. It is above all on such material that I have relied in the psychological section of my Philosophy of Religion, in that I at the same time compared the types I found in this way with the types presented by the great popular religions, especially their founders. But it will naturally be of importance to obtain a broader basis, and this could be attained if information and utterances could be collected from larger circles, also from non-literary circles. Attempts to procure such meet with many \opage{43}difficulties, partly because of the helplessness that will prevent many from expressing themselves clearly, partly because of the shyness that will prevent many from speaking about what is most inward in their lives; the mere question will perhaps even at times be regarded as a profanation. But since psychological research, like all other research, is pursued for the sake of truth and clarity, it will be fully justified to go this way and try how far it can lead. It is especially in North America that this has been attempted. In North America the religious life stirs from the practical side, and the prevailing Methodism with the revival movement attached to it has naturally directed attention to the mental processes that go on during religious experiences. It is also the psychology of conversion that the American psychologists have especially occupied themselves with. The leaders of the revival movement had already begun a statistics of the frequency of conversions, and the first attempts from the purely psychological side were in part made in cooperation with them. By means of circulated question lists an attempt has then been made to obtain information about the states before, during and after the religious breakthrough, as well as about other types of religious development that did not have the character of breakthrough, about crises and doubt, etc. Through analysis of the utterances obtained an attempt has then been made to arrive at a description of certain main forms of religious development --- naturally within the given historical presuppositions (namely American Methodism). Very large numbers could not be in question here, and therefore not a strict statistical method. One has cast the net out at a single place in the sea of religious phenomena and described the species one found in one's haul. There has, as the following will show, really been gained new material for comparative psychology in this way.

The first in the series of these investigators is JAMES LEUBA with his treatise \textit{Psychology of religious} \textit{phenomena} (American Journal of Psychology. April 1896). From his childhood Leuba had moved in circles where religious breakthroughs were frequent, and later he has, from psychological interest, extended his experience by conversations with people who had experienced such, by attending meetings and by sending out question lists \opage{44}concerning the circumstances before, during and after conversion. ``The question lists,'' he says, ``were sent to persons who thought they had `experienced religion,' to ministers and to leaders of the Inner Mission, and through them also to a large class that it had not been possible to reach directly. They were moreover announced in `The Presbyterian' in Philadelphia, in `The Christian' in London and in `The Outlook' . . . A small percentage of the answers received contained the phrases that are current at revivals and in mission houses and smacked so decidedly of religious affectation (religious cant) and so little of personal experience that they had to be set aside; but most of the answers bore inner testimony of serious endeavor to give an exact description of an experience to which immense importance was attached . . . The yield of our questions was greatly limited by the lack of confidence of certain people in the usefulness of applying psychology to religious phenomena, by the difficulties which ordinary people, who are not practiced in introspection, feel in having to give an account of their own experience, by the hostile attitude of some ecclesiastical authorities, and finally by the aversion that the invitation to investigate the nature of religious experiences sometimes aroused.'' Regarding this last point Leuba writes in a later treatise as follows: ``The ground tone of deep, sometimes even tragic seriousness that runs through most of my documents would move even an ironic skeptic to awe. The real value that is attributed to religion reveals itself when the intellectual assumptions on which the religious life is thought to depend are threatened . . . A lady of my acquaintance, to whom I gave a question list, recoiled in great horror when she had read the first questions. Her unclear grounds for her refusal to answer showed that she did not dare to look her faith straight in the face; the mere thought of it struck terror in her. Another lady, who with striking vehemence had declined a personal request for an answer, wrote a year later of her own accord to defend her conduct. Her faith was, at the time when the question came, in a precarious state, and she fought with the energy of despair that her heaven should not lose its God. Her chief means of defense was to drive from her consciousness \opage{45}anything whatever that could threaten her faith. The critical investigation I proposed wounded her to the quick. These moving examples of a class of our fellow men who prefer to stick their heads in a bush rather than let the hard facts rob their world of the soul that gives it grace, awaken to significant reflection on the connection between faith and need.''

In a series of later treatises Leuba has continued his investigations \textit{(The personifying passion in youth.} The Monist. July 1900. --- \textit{Introduction to a psychological study of religion.} ibid. January 1901. --- \textit{The content of religious consciousness.} ibid. July 1901. --- \textit{Les tendances religieuses chez les mystiques chrétiens.} Revue philosophique. Juillet—Novembre 1902). What characterizes Leuba is his interest in the individual cases in their multiplicity and peculiarity. He has brought together a collection of religious self-characterizations, which show how many nuances and transitions there are in this field, and how little our usual rubrics suffice. ---

Related in design and in significance to Leuba's investigations is a recently published treatise by LUCIEN ARRÉAT: \textit{Religion in France.} (The Monist. January 1903). In his investigation of French Catholicism he gives an interesting counterpart to Leuba's investigation of American Methodism. ---

EDWIN STARBUCK \textit{(The psychology of religion.} An empirical study \textit{of} the growth of religious consciousness. London 1892 disposes of a very extensive statistical material, which he has subjected to a thoroughgoing analysis. He conceives conversion as a process through which the young person is to make what is socially handed down into full seriousness and reality in himself, and therefore as a sociological phenomenon that can be demonstrated among all peoples and in all mental domains. Of particular interest is his demonstration of different kinds of conversion and his comparison of the process of conversion with other forms of religious development.

GEORGE COE (\textit{The spiritual life.} Studies in the science of religion. New-York 1900) has by a similar path as Leuba and Starbuck especially investigated the influence of temperament on \opage{46}the character of the religious development and its relation to other co-operating elements.

GEORGE SANTAYANA has in his work \textit{Interpretations of Poetry and Religion} (London 1900) especially brought forward the poetic element in religion and has thereby been led to place himself in opposition both to orthodoxy and to the current criticism. For him religion is the highest poetry, while he at the same time seeks to show how poetry at its height passes over into religion. ---

On the whole the standpoint of the American psychologists toward the religious life is a sympathetic one. They energetically assert the independence of psychology in relation to all theology and speculation, but they immerse themselves with enthusiasm and interest in the religious phenomena. Their own religious standpoint does not seem to be markedly anticonfessional, even though they think they can demonstrate the necessity that the confessional limitation give way to the valuable multiplicity of paths and types that unbiased observation discovers. As far as possible I shall in the following discussion of their investigations also seek to determine to which type they themselves belong, in order thereby to determine the personal equation and indicate a limitation for the significance of their investigations.

\underafsnit{a. Breakthrough and Continuity in Religious Development.}{a. Breakthrough and Continuity in Religious Development}

STARBUCK found in the material he collected two types of religious development. The one moved through a more or less sudden decision on to a new direction of life; the other proceeded more continuously.

\textit{The breakthrough type} again presents differences of interest

The age at which the breakthrough takes place can vary. It is most frequent in the young years and occurs somewhat earlier for women than for men; later than at the age of 40 it occurred rarely (within the 1265 cases where the age could be ascertained). --- The occasion for the breakthrough varied. A large part of the cases was due to revival meetings (revivals). Some took place during ordinary church service. \opage{47}In others again no particular external occasion could be demonstrated. --- As regards the motives, fear occurred far more frequently than hope, and altruistic motives (ideality and love of humanity) were rare. It is characteristic that the feeling of sin was rarer at revivals than on other occasions; Starbuck thinks that at revivals one does not get time to immerse oneself in oneself and brood over one's conduct and one's past.

Most important are the differences with regard to the character of the breakthrough itself. Two forms appear here, which, however, can be demonstrated in their purity only approximately. In the one form it is a habit that is broken, a vicious state from which one escapes; a deep feeling of sin then precedes. In the other form it is a new light that breaks in; a lack is filled and a positive fulfillment is reached; a feeling of incompleteness then precedes. The former form has a negative, the latter a positive character. The negative form requires complete overcoming and repression of the past in order to come about; the positive form requires only a small impulse to take effect. Starbuck remarks that within the limits of his material the positive type (spiritual illumination, preceded by a sense of incompleteness) is more frequent than the negative (escape from sin, preceded by a feeling of wretchedness), and he stresses the questionable character of the tendency prevailing in the Church to make the negative type the one right and normal one. Those who struggle for fuller light must be treated differently from those who struggle to be freed from darkness. In the material collected by Leuba too the positive type occurs most frequently. Starbuck and Leuba do not discuss the cause of this. It is not quite certain that the positive type is in reality the most frequent; perhaps the negative type is more shy and disinclined to express itself, so that therein lies the reason why fewer cases of it can be ascertained.

Of very great interest is the interaction of the conscious and the unconscious in the breakthrough. There are many different degrees of consciousness here, which probably is connected with there being different degrees in the inner hindrances that are to be overcome. A presupposition for the breakthrough \opage{48}is --- within both forms --- that a growth has taken place beneath the threshold of consciousness (within the sub-liminal consciousness). Now it is not given that the direction of the conscious striving coincides with the direction of the unconscious growth, --- and since the assent of the will to what is unconsciously laid down is nevertheless necessary for the breakthrough to occur, it is no wonder that so many psychological riddles present themselves here. What is to be brought about is, after all, that what has hitherto stood only as object, as a distant and exalted ideal, is to become that very thing which lives and stirs in the individual. What has hitherto been only object is to become subject.

Once the breakthrough has occurred, a feeling of newness arises, as a rule also of joy and peace, and frequently of self-forgetfulness and love of humanity. What has hitherto held the self back within narrow barriers is displaced, and an expansion can take place.

Psychologically viewed, conversion in both its forms is not a unique phenomenon. There is a series of analogous experiences, as when uncertainty and dejection are suddenly replaced by cheerful decision, or when a long effort to acquire a skill is as with one stroke replaced by mastery over the difficulties, or when a state of feeling announces itself suddenly through contrast or prepared by inner organic processes. Such phenomena occur especially in the age when the young person is to live himself into the habits and traditions of the race and make them his own --- is to expand his self to embrace larger interests, break down the barrier between the individual and the social will. Physiologically viewed, this task is set at a time when growth is approaching its completion, and when new tasks and functions begin to demand their due. It is a time of movement and of oppositions; both in the organism and in consciousness the opposition between the old and the new emerges sharply. But through religious influence this opposition is often strengthened in an unnatural and unhealthy way. One does not always recognize, as already remarked, the difference between the two types, but treats them alike. Starbuck cites an example of how it can lead to despair \opage{49}when all are to be forced into one definite scheme of development and are compelled to make experiences that are impossible for them. ---

In opposition to the breakthrough type in its two forms Starbuck sets \textit{the continuity type}. There is here no consciousness of decisive change, of opposition between old and new. An involuntary awakening takes place, which indeed has a certain likeness to the positive form of breakthrough, but lacks the experience of sudden change. Development proceeds continuously from childhood through the young years. Yet a period of ``storm and distress'' (storm and stress) can come, which is connected with the organic changes of the years of transition and can bring anxiety and tension. It is the period when the young person is to live himself into a larger world and make the acquisitions of the race his own. Manifold demands and new possibilities present themselves. In intellectual form (and somewhat later) the unrest thus awakened appears as doubt. The point is to hold fast with clearer consciousness to what one had hitherto accepted on authority and with childlike credulity. Such a period was very frequent among those whom Starbuck came to know more closely; it could begin, run its course and conclude without abrupt transitions. Often it led over to a time when the individual stood estranged toward what had been handed down and regarded it as wholly worthless. This period of estrangement (period of alienation) was often due to relaxation after the unrest or doubt; but it could also have its ground in the individual's now wishing to assert himself against the forms society wished to lead him into, or in delight in the unrestrained use of faculties and powers. In place of the religious tradition aesthetic, intellectual or other interests often step into the foreground in such a period. At the end of the period of youth a reconstruction as a rule sets in: the individual has won a definite faith, which does not always agree with the received one, but as a rule is a mixture of independent production and resumption of the old. Starbuck is astonished to find this period so definitely marked. He explains it by life's circumstances demanding a definite standpoint in the practical respect, so that perhaps even in the race a tendency to such reconstruction after the \opage{50}agitated years of youth has developed. --- Through all these periods development within this type proceeds so continuously that it is often only through comparison of points lying far from one another that one becomes aware of them.

Characteristic of the conclusion of development is --- and this holds for both main types --- a certain change in the religious ideas and in the whole conception of religion. Even though the confessional ideas are essentially retained, a pervasive shift can nevertheless as a rule be demonstrated from the outer to the inner, from the dogmatic to the ethical, from the isolated personal to the consciousness of participation in a more comprehensive life. Religion is conceived more and more as an inner life that is part of a universal order of things.

\textit{If we now compare in conclusion the breakthrough type with the continuity type,} the former appears, as is natural, with far greater unrest and with ``storm and distress'' to a more marked degree than the latter. The breakthrough type is more receptive to outward influence; the opposition between the old and the new asserts itself more strongly in it, and the life of feeling is more vehement. The breakthrough type moreover adheres far more than the continuity type to the definite confessional forms, to the organized institutions. Its inner experiences it prefers to express in biblical and confessional terms and has no need to form independent forms and expressions or find new analogies. Such a need, on the other hand, is very frequent within the continuity type, where also the conception of religion as an inner life occurs more frequently. --- Starbuck's general result is that the development of the religious life under confessional forms can be illuminated from general psychological and physiological points of view. Dispositions and impulses of general human nature unfold in religion, and also the different forms in which the religious development takes place are known from other mental domains. He himself regards religion ``in its highest form'' as a radiation, an interplay and an idealization of tendencies that are already present beforehand in human nature. The forms and expressions the race has found for this are what the young person is seriously led into in the years of transition, and the course and \opage{51}crises of development are conditioned by the difficulties to be overcome here. He is inclined to think that the inconveniences that appear in the course of this development are due far more to the older generation's lack of understanding of the demands of human nature than to oddities and faults in the young. --- With this last remark may be compared a statement by Professor PEABODY in his treatise \textit{The religion of a college student.} (Forum. June 1901. p. 447): ``For one young man who is removed from religion through the influence of science, philosophy or art, there are quite certainly ten who go thus because of unwise or unpractical religious instruction . . . . It is not the University that is to renew the agreement with the Church; it is the Church that must open its eyes to the wonderful expansion of the intellectual horizon which every student nowadays has access to experience.''

\underafsnit{b. The Influence of Temperament.}{b. The Influence of Temperament}

GEORGE COE too stresses the importance of a psychological study of religion in its present state. And it is also the period of youth to which he has especially directed his attention. This period is, he maintains, far less known and understood than childhood and mature age. What takes place in this period is a transition from being taken up with the particular to grasping the universal, from being overwhelmed by the outer to immersing oneself in the inner, from letting oneself be led by authority to being determined by reasons, and from stopping at the given to striving toward the ideal. Questions are asked more searchingly about the significance of life, and the social interests begin to awaken. All this brings it about that a need is felt for a more personal appropriation of the content of religion than there is possibility for in childhood. Conversion or a corresponding process of inwardization is therefore a natural link in the development of the period of youth, and Coe also in fact finds that such a change is extraordinarily frequent (especially at the age of 16). He recalls that among peoples around the earth it is common to place a religious initiation at this time. Often the young man is sent out into solitude to wait for a revelation from \opage{52}the god of the tribe and by its help to become clear about his calling. In agreement with Starbuck he maintains that among young Americans within Christianity it is by no means merely fear or remorse that are motives in religious development, but also more indeterminate states: unrest, emptiness, feeling of incompleteness, longing toward the infinite, --- even though most express their experiences in the confessional language. It further appears throughout that purely intellectual problems and motives are subordinate; the reasons are found afterward; one seeks first to live the life that expresses the deepest self, and afterward one perhaps looks to see how one can by the help of reason show its justification. And it is not any single experience that leads to the goal. Most of what youth demands decisively all at once, one can only live oneself into gradually, just as life unfolds with its work, its disappointments and its victories.

Coe wishes to study these phenomena without theological and metaphysical prejudices. The psychology of temptation, e.g., is and remains the same, however one may explain to oneself the origin of the impulses that carry it with them, --- whether one thinks that they derive from the Devil, or that they are remnants of the animality from which humanity has developed. The same holds of the psychology of conversion, with which Coe has especially occupied himself. It is his task to point out the most important factors which co-operate in this psychical phenomenon. He goes here the way of making an exact investigation of a relatively small series of individual cases. Besides sending out question lists he used detailed conversation with individuals, procured others' testimony about them, used hypnosis, sought by all possible ways to arrive at an individual characterization. The number of individuals he investigated was 77 (52 men, 25 women); they were almost all students, the men on average 24 years old, the women 22 years; most had been brought up in the Methodist church, which indeed lays so great weight on personal religious development.

It then turned out that in such young people great differences appear with regard to conversion. \opage{53}In this respect three classes of individuals can be distinguished: A) those who expected a transformation (transformation) and really experienced it; B) those who expected a transformation, but did not experience it; C) those who were not satisfied with what they experienced, or who nevertheless strove to attain a fuller experience.

On closer investigation it now appears that class A were predominantly people of feeling (sensibility distinctly predominant), whereas class B were predominantly people of intellect (intellect distinctly predominant). And whereas class A most closely recalled the so-called melancholic and sanguine temperaments (the slow-intense and the prompt-weak varieties), class B recalled more the so-called choleric temperament (the prompt-intense variety). Class C presented differences that recalled now A and now B. Temperament thus proves to be an essential factor.

With this there now combines another observation. In class B there occurred next to no automatic phenomena in connection with the religious states, neither in the domain of the life of ideas (strong dreams, hallucinations), nor in the domain of the organs of movement (irrepressible laughter, starting, clapping of hands, talking aloud to oneself, etc.), whereas such phenomena occurred far more frequently in class A (and in C). --- When the classes were moreover examined with regard to susceptibility to suggestion and hypnosis, it turned out that A's suggestibility is predominantly passive, in that the individuals here quite conform to the prompting given and follow it in detail, whereas class B's suggestibility can be called spontaneous, in that the individuals either give themselves their promptings or independently add something to the original prompting.

Besides the aroused expectation, temperament, tendency to automatism and suggestibility thus play a role in conversion. The character of the religious experience is determined by these factors. ---

Like Starbuck, Coe stresses very strongly the danger of regarding a single type as alone justified. If an expectation is aroused of experiencing a more or less sudden \opage{54}upheaval of mind, and the fulfillment of this expectation is, according to the individual's whole psychological nature, impossible, then it can lead to despair. Coe thinks that really it is only the melancholic and the sanguine temperament that the Church has recognized in its worship, its hymns and its morals\textsuperscript{11}). In the type of Christ that prevails in the Church this shows itself especially; it is the Man of Sorrows, the suffering sacrificial lamb. The real Jesus suffered because in spite of all opposition he held fast his thought to the utmost; he was an active, not a passive individuality; there was a choleric element present, which has vanished from the traditional image. --- The Church must be as broad as human nature is. But it has accustomed itself to look at things with feminine eyes: therefore it is especially the men who turn away from it.

Coe finds, in spite of his evident sympathy for religion, no scruple about a thoroughgoing psychological discussion of the religious states. If we --- he says --- want to understand religious experiences, we must not explain them by God's unsearchable ways, but by demonstrable peculiarities in the nature of man; we are not to speculate about divine grace, but to study the hidden life of the human mind. Like other philosophers of religion in our day he finds the final measure for religious values in the ethical. What matters is what is gained by religious development for a deeper, more richly contentful active life.

\underafsnit{c. Differences in the Relation to the Confession.}{c. Differences in the Relation to the Confession}

There are here, as already mentioned, two investigations, one by Leuba on American Methodism, another by Arréat on French Catholicism.

LEUBA has studied partly the religious life as it appears among contemporaries (in his native country North America), partly the religiosity of the mystics. Both studies lead him to make a definite distinction between the standpoint of the historian of religion and that of the psychologist of religion. The historian's subject is not so much the religious experiences themselves as the products of these experiences present in worship, myth and dogma, as they \opage{55}take shape within a society. To arrive at the immediate religious experience one must turn to the individual religious consciousness. By it all religion is in the end borne and conditioned, and in the higher popular religions this relation becomes especially clear. It is not the gods but the relation to the gods that interests the psychology of religion. Its task is to give a religious dynamics, i.e. a doctrine of the motives of the religious life. (The expression ``religious dynamics'' had been used by Coe before Leuba). To obtain material for this from contemporaries Leuba has by sent-out question lists obtained utterances from a series of persons who stood at very different standpoints, but all within the Anglo-Protestant world. The questions concerned the grounds of their participation or non-participation in worship (public or private worship), --- the circumstances and localities that exerted a special religious influence on them, --- the most characteristic religious experiences they knew, --- the differences they made between religious and non-religious states and actions. He has succeeded in collecting a series of characteristic utterances. Statistics is not in question; but Leuba thinks he can demonstrate certain peculiarities of the religious life in the circle he has to do with.

A main point is for him the result that what expresses itself in religion is not a striving for knowledge, but a striving for being. Religion springs, in his view, from a feeling of fragmentariness (unwholeness) in connection with a feeling of inner division and disruption. The former feeling brings with it a need to become a partaker in a larger, more comprehensive life, whereas the feeling of division brings with it a need to be led by a single end (unification of motives). The feeling of sin is a form of this latter feeling. To the two motives mentioned are joined further altruism and the feeling of dependence: the first is related to the need for unity, the second is conditioned by the inner resistance that is to be overcome in efforts to attain unity. The great founders of religions paid no heed to theoretical problems, but strove after ideal perfection and inner unity. And there manifests itself nowadays \opage{56}more and more a need to go back to the immediate religious experiences from which religion has always from the beginning grown. Although one has not yet had courage or intellect enough to do this fully, still the first steps in this direction (to return to plain religious empiricism) have been taken, especially within Protestantism. In silence there is waged --- alongside the struggle for existence between individuals --- a purely inner struggle, which in severity and tragic seriousness often surpasses the outer. It is waged in different ways within the different religions and finds expression in different forms: but always there is bound up with it the need for a greater and at the same time more undivided life. The religious need is really no quite special need, but has its deep and lasting psychological and physiological foundation.

In comparing religion as it appears in more recent times in contrast to earlier forms of it, Leuba finds a series of differences.

In primitive societies religion is a composite product of all of man's fundamental wants and strivings, both theoretical and practical. The theoretical motives have now separated out and find their satisfaction in science; it is not on them that the significance of religion in more recent times rests, and already in antiquity Buddhism stands as testimony that these motives are not the decisive ones. Even if all intellectual elements hitherto attached to religion fall away, the psychological essence of religion, as it has been described above, will be able to persist.

A further significant difference between more primitive and more modern religion Leuba finds in this, that whereas formerly the feelings of obedience and awe were predominant, these are now --- in any case in the Anglo-Saxon world --- not so frequent as is assumed. Not the obedience that the catechism inculcates, but the clear understanding of what is right and good, is now the motive to action, and the feeling of unity and freedom is displacing submission. Connected with this is the feature that religious truth is no longer regarded as something that can and ought to be grounded in an objective way; the \opage{57}decisive thing is what satisfies the personal need and yearning. Many of those questioned expressed this with the greatest candor. The motives of feeling had decidedly the preponderance over all intellectual interests, and there was perhaps even complete indifference toward the purely objective significance of the religious ideas, even in individuals who otherwise in more specialized fields showed themselves in possession of critical sense and capacity. One does not \textit{know} God; one does not \textit{understand} him; but one \textit{uses} him. Not God, but life, more life, higher, richer, more satisfying life is the end of religion. In the past this shows itself clearly in the so-called mystics, and here Leuba's studies of mysticism connect with his studies of contemporary phenomena. The mystical life must not be interpreted and judged by the concepts and dogmas the mystics employ. Its peculiarity is only a more extreme and concentrated form of the peculiarity of all religion: a yearning for higher and more complete satisfaction, for peace and unity in the will. The mystics' yearning is more passionate than that of ordinary religiosity, and the concentration upon a single goal is stronger. Goal and direction are the same, but they proceed with far greater energy.

Yet another characteristic feature of the religiosity of more recent times Leuba finds in this, that so many and considerable individual differences appear, even within one and the same civilization. Under the influence of the infinitely larger horizon life has acquired since the time of the Councils, religion too has undergone changes, and differences show themselves not only in the outer things about religion, but in the individuals' innermost mental life. The wants are different, and the hopes are different. And yet in theology one proceeds from the assumption that one and the same food suits all. ---

Within his collection of self-confessions Leuba demonstrates a whole series of nuances, from a standpoint where the religious and the confessional coincide to such a degree that the individual is not able to express himself about his inner experiences except in biblical and confessional terms, to such standpoints where no meaning whatever can be connected with that kind of \opage{58}expression. Between both extremes lie the most interesting cases; for neither from dogmatic position nor from dogmatic negation is there any considerable psychology to be learned. Only a few features can be cited here from the interesting self-characterizations Leuba has collected.

A scientifically educated young woman, adherent of the Presbyterian church, distinguishes definitely between what is due to her confessional upbringing and what is due to the need for something that has eternal value (a desire for things of permanent value). These two things do not exclude each other, but neither do they coincide. When she takes part in the worship of her church, it is partly from habit, but partly also because she needs to enter into relation with a higher being and to be freed from selfish thoughts. The difference between religious and non-religious thoughts and feelings is for her a difference of degree and rests on the fact that the former appear as more significant and stronger. But the ideal need is never fully satisfied. ``Whereas the current religious assertions declare religion to be the only thing in the world that can satisfy,'' she says, ``for me its value rests on its constantly driving me forward to new exertions.'' --- Other correspondents emphasize love of humanity, childlike mind or gratitude as the essential in what they call religion. --- A physician, who takes part in the worship of the Church, although his own private devotion has for him a more decidedly religious character, states among other things the following: ``I am far from having the anthropomorphic idea of God that I had in my young years. By God I understand the power or the principle that brings about all development in the universe toward a higher order of things, and which by its laws raises the stages of life from the unconscious to the conscious, from the animal to the rational, from the perishable to the imperishable. In myself I call divine all drives, tendencies and exertions that lead to developing in me greater receptivity to beauty, truth and justice, higher reason and capacity for moral progress''. --- A related statement from another correspondent is the following: ``I have given up every attempt to imagine God and no longer feel any need for an image \opage{59}or a clothing of the divine principle which it is a necessity for me to believe in. My religious relation to what one calls God I can describe from two points of view: 1) I feel in my innermost consciousness a need to be One with a great principle, a need for unity with Something that I cannot define, but which surrounds me, in which I am so to speak immersed, and with which I wish to be in harmony, an all-embracing Present . . . . 2) I have almost uninterruptedly a deep consciousness of standing far below an inner ideal, a type that is laid down in me and that I constantly have before my eyes. I feel that I \textit{must} realize this type, which constantly dwells in me and seeks to lead me. Experience has taught me that I become happy by following it, unhappy by moving away from it. To these two fundamental instincts or intuitions I have reduced all my religious ideas. I am a disciple of Christ and strive to imitate him; my ideal is to become One with the Divine, as he was. But there are many of the assumptions of the great Christian communities that I do not share.'' --- A woman who lives with husband and child in the country in constant struggle with poverty declares that for her religion is truth, and by truth she understands the law that holds for our being. ``When I say God, I mean the Good, the law of my being . . . I believe that life is development, growth . . . and that we have as much heaven and as much hell here in life as we ourselves will . . . The Good is our ideal; the Evil is our teacher, who is to make us understand the Good . . . Human beings cannot be divided into good and evil, but into developed and undeveloped . . . I know no religious act except life itself . . . Perhaps I was nearest to religious experience when it first became clear to me that I was free, free from faith, doctrine, form and sentimentality, free to grow, to become All that lies in my nature . . . If you cannot find any religion in this, you will hardly find any religion in me . . . . I am against the doctrine of the imitation of Christ, because it conflicts with independent thinking, and because mental and spiritual development rests on independent thinking. We might just as well let another eat for us as think for us. Enthusiasm one \opage{60}mind can inspire in another, but not Existence . . . . Christianity has been a great factor in the development of the human mind, but the Christianity of dogmas one has now almost outgrown . . . . We are not about to lose our religion, but we are about to find ourselves. We conquer fear and learn to have trust in the divine law that rules over life.''

All these examples belong to the group that thinks itself positively religious, although it is often, especially as regards the last example, doubtful with what right. But they are characteristic expressions of the need to give an account personally of what one --- the single individual --- can find of value for his life in the given positive religion. Through some transitional cases of more indeterminate form we come in Leuba's series over to decidedly non-religious forms, which --- however far one may stretch the concept of religion --- cannot come under it. Partly there is lacking here religious experience in any sense of the word, partly there are definite intellectual and moral grounds that lead away from religion --- admittedly because religion and the Church system are assumed to coincide. Thus one correspondent remarks: ``I deny the moral superiority of the churches. They never take the lead of a moral reform. They were on the wrong side in the slavery question. They are cowardly, petty and servile in the great industrial questions of our time.'' --- Leuba thinks that the number of non-religious persons is far greater than is commonly believed. When there are not many utterances from such persons in his collection, it is, in his opinion, partly because his question lists have not reached them, partly because they are not always clear about their real standpoint. ---

Leuba has not himself concluded the analysis of his material; his treatises available so far are only results of a first treatment, and he ends really with an unanswered question. In the fundamental religious motive, the religious tendencies he finds nothing that is peculiar to religion, but particular forms of the need for self-preservation and self-development. This conception he finds decidedly confirmed by his correspondents' utterances. He rejects every conception of religion \opage{61}according to which it is to aim at knowledge, whether this knowledge is to give a revelation or to establish a mystery. (The Monist. July 1901. p. 571---573). But at the same time as he in his latest treatise (Revue philos. Novbr. 1902. p. 487) refers back to these statements, he states: ``What then is the element that distinguishes the religious life from the secular (for there is, after all, a difference between them)? It is the one means or the several means that man employs to attain his wishes, and not these wishes themselves, that separate religion from the rest of life. Psychologically expressed: it is not the will (in Schopenhauer's meaning of the word), but knowledge, that separates them. When one realizes that the difference between the religious people and those who are not (insofar as there are such) does not consist in striving, wish and tendency, thus not in the very sources of life, but only in the means employed to reach the same goal, religion appears in another light than that in which it ordinarily stands in people's eyes.'' --- A closer explanation of these hints will hopefully be brought by Leuba's following works. But when he definitely rejects the intellectualist conception that it should be the \textit{end} of religion to give knowledge of any kind whatever, and yet thinks that the difference between religion and other spiritual forms of life \textit{rests} \textit{on} knowledge, then the meaning can well only be that the difference rests on the way in which knowledge is used in religion as \textit{means}. And means must here mean means of interpretation, that is, symbol. I also believe that Leuba would be able to draw from his material a result in this direction. There is in all religion a need to form what is experienced in thoughts. In confessional religion this need is satisfied immediately by the received ideas; in the transitional forms other expressions are sought to interpret the significance of these received ideas themselves for life. At the base there always lies an experience of the relation between ideal and reality, an experience that --- wherever religion at all arises --- passes over into faith and trust toward powers that are assumed to work in the direction of a harmonious relation between ideal and reality. I find my hypothesis on the essence of religion as faith in the \opage{62}persistence of values confirmed by the investigations of Leuba and his colleagues.

On the other hand it seems to me that Leuba lays too little weight on the influence of prevailing or received ideas. It is fateful for the whole character of a feeling to become attached to a definite circle of ideas; it thereby acquires a direction and a content that decide its quality, and when the connection between it and the ideas it has hitherto been attached to is to be torn apart, it can happen, or at any rate threaten, that the feeling disappears at the same stroke. One can indeed with Leuba distinguish between ends, which can be said to be common to religion and other spiritual forms of life, and means, which are different for the former and for the latter; but it must then surely be maintained that the very fact that an end can be attained by different means indicates a difference in the end itself. There is no wholly external connection between end and means in such domains as this concerns. The sting of the religious problem lies precisely at this point. History teaches, after all, also how hard human beings find it to put up with differences of quality within religion, and the dispute over the limits of the very concept of religion rests likewise on the question whether through a difference of quality we get outside religion's whole modality (to use the terminology of the psychology of the senses). It becomes thereby of particular significance that there are natures that are able to find expression for their inner life only in \textit{common} symbols, and who, if they are compelled to give them up, feel wholly homeless.

Yet it must be repeated that Leuba has not yet completed the working up of his material. It will be of interest to hear the words with which he concludes his preliminary communications. ``Whether the differences between religion and non-religion are only superficial, or whether they are fundamental, and in what degree, --- whether upbringing in the home and other circumstances in and of themselves can explain these differences, or whether they have deeper roots in particular innate dispositions, --- whether the one is destined to displace the other or is merely a preparatory stage for it or not, --- it is a question \opage{63}that must be passed over in silence in this preliminary treatise. Only this I wish to emphasize here: even between individuals of the same society and of about the same intelligence and culture there are differences that are sufficient to make the one a human being for whom religion is the All of life, and to leave the other standing wholly indifferent and inaccessible to it, even if it be interpreted in the widest sense.'' ---

LUCIEN ARRÉAT has --- by similar paths as the American psychologists --- sought to find different types within French Catholicism. One notices at once that the circumstances are different. The difference between Catholicism and Protestantism is undoubtedly connected with racial differences. Whereas from another side (Charles Renouvier) it has been regarded as the fall of the French people that it did not join the Reformation, Arréat maintains that this conflicted with the peculiarity of the French race. Protestantism offended the sense of beauty by its destruction of images and buildings, and what it robbed the imagination of it gave no compensation for in the domain of thought. Protestant individualism demands of the single individual an independent discussion and understanding, which conflicts with the Latin race's disinclination to theological discussions. The Catholic sees the doctrine as fixed once for all, and his inwardness breaks forth in the personal relation of feeling, the enthusiastic adherence to the fixed eternal truth which the Church guarantees. Arréat has therefore also found in his material no hint of the transition from feeling of obedience to feeling of freedom in the relation to God, which Leuba thought to find in Anglo-Protestantism; Catholic religiosity is essentially submission. --- According to his observations Arréat describes four different classes within French Catholicism.

The first class, which is found especially in marked form in southern France, comprises the largest number of Catholics and has representatives in all classes of society. Piety consists here in observance of the religious customs, which are conceived in a way that approaches idolatry, which shows itself especially in the strong flocking to the places where \opage{64}miracles occur. Here the racial character is at work whose inclination to superstition Caesar already mentioned.

The true Catholics make up the second class. Through participation in the means of grace of the Church they strengthen the feeling of security of being on board St Peter's ship, which in the midst of the unrest of the times is steered by an infallible helmsman. The Catholic religion is praised for its sublimity and its captivating beauty, for its ideal of perfection and for its active love of humanity. Without need of mysticism or of independent thinking there expresses itself here a mental equilibrium, which is nevertheless sometimes connected with an express conviction that it is only the mighty authority of the Church that keeps doubt and danger away.

To a third class are reckoned those who support their ecclesiastical faith by more or less conscious thinking. Many clergy belong here, and to this class Arréat thinks he can also reckon many who declined to express themselves. One protests here against religion's being a matter of feeling: it grounds itself on rational knowledge, although grace brings the final decision. The Catholic doctrine stands as a whole from which nothing can be separated, and which immediately captivates the soul. Yet it is sometimes said that the Church is only the outer and perishable garment of faith, which will be able to be transformed into new forms; the properly religious is placed in the need for the infinite. By such statements the transition to the last class is made.

In many men and women who declare themselves good Catholics a certain state of dissolution can be traced. They choose, e.g., among the Church's doctrines and customs and hold to what they have use for. Often it is only an express act of will that holds them back in the Church; one stresses the social motives, the opportunity to practice love of humanity, e.g., as that which binds to the Church. Or one conceives the religious as ideal striving, pays homage to a practical idealism without metaphysics and uses the Church's worship only as a means of mood. One expresses the most inward sympathy with the believers, because they do not, like many free-thinkers, give up the sublime poetry of religion or are blasé toward the mighty \opage{65}problems to which the word God points. A last step is taken when one goes over to Positivism or by another path seeks a wholly new expression for one's view of life. This is, however, unusual. Either one returns to the faith of childhood after a time of unrest, or one sinks in spiritual respect to a lower level.

It is --- Arréat says in conclusion --- impossible at the moment to say whither the present crisis in the religious domain will lead, or what energy the religious feeling will be able to possess when it is not supported by faith in dogma or poetic myths: ``I see what one wishes to eradicate from modern mental life; I do not see what is offered for its enrichment.''

\underafsnit{d. Poetry and Religion.}{d. Poetry and Religion}

GEORGE SANTAYANA, in his work mentioned above, has discussed the relation between poetry, religion and science in a way that gives him a place in the ranks of the psychologists of religion. For him it is the imagination, the image-forming faculty, that underlies all those three spiritual manifestations of life. The images that we can use in practice, and with the help of which we can predict future experiences, we call ideas of the intellect, whereas images that cannot be exploited in this way continue to be called ideas of fantasy. Intellect is a fantasy that makes empirical science possible. Under the influence of experience we choose among the images that arise involuntarily. But among the fantasy images that cannot pass over into ideas of the intellect, two kinds must in turn be distinguished. When fantasy images, although they do not express empirical reality, nevertheless stand as expressions of life values, ideals and ends, they receive the character of religion, whereas they stand as purely poetic forms when they have no such significance for life. Religion and poetry are in their innermost essence akin, indeed identical, and differ merely in their different relations to life. Poetry at its height passes over into religion when it expresses the value of our experiences and gives great, guiding and inspiring anticipations. Conversely, religion passes over into poetry when it gives up the presupposition of containing metaphysical truths. \opage{66}What matters, in order rightly to understand this, is to distinguish between the world of values and the world of facts: value rests not on the given itself but on its significance.

Our will is infinite, but our experience is limited. Poetry leads us beyond this limitation and thereby makes possible a higher development and exercise of the will, in that its highest end can now appear with greater clarity through poetry's help than mere experience allows. A striving, a passion we understand rightly only when we see what is sublime in its whole and full unfolding. But because poetry becomes religion, it does not cease to be poetry. On this point not only the orthodox but their critics as well are guilty of error. Already the Greek philosophy's criticism of polytheism was unjust, in that it took no account of mythology's close connection with life. A great religion is a product of the human spirit in its entirety, and in relation to that, criticism, which is guided by the special, intellectual interest, stands as limited and piecemeal. In Christianity historical idealism took the place of the naturalism of the ancient religion; the harmony that the old religion had sought to represent in space and in nature, the new sought to represent in time and in history. Christianity obtained its foothold in history by having its root in a national faith that had taken shape under the trials of the Jewish people; but out of this it formed a great poetic expression of experiences that human beings knew from the innermost depths of their souls. It was the great autobiography of humanity that was here made the center of the life of the world. The significance that a religion has and retains is not necessarily the one that it had for its founders; they are so captivated by their visions that they regard the content of these as external facts. Religion is not different from superstition in psychological respect; the difference rests not on origin but on value. Even if ethical value was not the ground of the coming into being of religions, it is nevertheless the ultimate standard of their significance. ---

This whole line of thought, which Santayana carries out in detail through interesting characterizations of religious and poetic works, acquires a special interest when one compares it with \opage{67}some of the observations of Arréat, Leuba and Coe, and when one at the same time recalls the idealistic or mystical tendencies that are stirring in the literature of the most recent time. What I miss in Santayana is a more penetrating psychological analysis of the relation between experience and faith; but otherwise it has been of great interest to me to find a way of looking at things akin to my own in a work that appeared at about the same time as my own attempt in the philosophy of religion. For comparison I permit myself to quote a single statement from my book (Religionsfilosofi p. 334 f.): ``Perhaps it is the case that poetry is a more perfect expression of the Highest than any scientific concept could be. By poetry, then, one must not here think of loose moods and fantasies, but of the involuntary and vivid form in which what is really experienced clothes itself in strongly moved moments . . . It is a poetry that neither stands in opposition to will nor to thought, although it comes especially into being where thought and will meet with barriers . . . . All high art presupposes such a poetry of life and is its conversion into clear forms.''

\underafsnit{e. Comparison between Different Series of Types.}{e. Comparison between Different Series of Types}

In my Philosophy of Religion I have, relying especially on a biographical and historical material, presented a series of typical differences in religious faith (p. 111—120) and, in another connection, a series of differences with respect to the original dispositions that condition the peculiarity of religious life in individuals (p. 251—256). If one compares these series with the types that the American psychologists of religion have set up, one will find above all one opposition, which plays a leading role in all the presentations, namely the one that I have described as an opposition between faith as concentration and rest after the balancing out of disharmonies and faith as the highest form of the unfolding of life. Corresponding to this, among the dispositions, is the opposition between disharmonious and expansive natures. In Leuba, Starbuck and Coe alike this is the most prominent opposition. But it also appeared that they lay quite special weight on the expansive type and are inclined to regard the \opage{68}disharmonious or discontinuous type as a special case of it. There is of course always some resistance to be overcome, some crises to be passed through, and it is only when resistance and crises become particularly strong that, in the opinion of these investigators, we pass over from the expansive to the disharmonious type. As to the reason why the expansive type occurs so frequently in their material, I have expressed a conjecture above. One might perhaps add that the authors mentioned (and the same evidently holds of Santayana) themselves personally decidedly seem to belong to the expansive type, in so far as one can come to know their psychical habitus from the way in which they perceive and evaluate those of others. For my own part I would be inclined to regard the difference between the two types as greater and more fundamental than those authors assume. Perhaps there lies in this opposition the germ of a spiritual struggle that again and again will come to stand between spiritual directions of life. In the deepest depths they ultimately do not understand each other. In Danish intellectual life we have a classic example of this opposition in the relation between Grundtvigianism and the Inner Mission.

Were I to reduce my series of types, this opposition would nevertheless remain. I have in all seven types, namely besides the two already mentioned [faith as 1) rest and as 2) unfolding]: 3) contemplation, 4) confident venturing, 5) resignation, 6) convulsive religiosity and 7) collier's faith (fides implicita). Of these, 3, 4 and 7 would fall under 2, whereas 5 and 6 would belong under 1. As regards the differences in dispositions, I set up, in the domain of the life of feeling and will, the following series of oppositions: between the Disharmonious and the Expansive, between the Passive and the Active, between the Affective and the Continuous and between the Idiopathic and the Sympathetic, --- and in the domain of knowledge the opposition between analytic and synthetic natures and between discursive and intuitive natures. It is easy to see that there is an analogy between all these relations of opposition, and that they can all be subsumed under the first, so that it is the opposition between the Disharmonious and the Expansive that, seen from this side too, comes to stand in the first rank.

\opage{69}In the conviction that in the future, to a far higher degree than hitherto, in understanding and judging religious phenomena, weight must be laid on individual differences, I am in complete agreement with my American colleagues.

\afsnit{4. Comparative Psychology of Religion.}{4. Comparative Psychology of Religion}

A comprehensive work in the field of the psychology of religion is to be found in the American philosopher WILLIAM JAMES's book \textit{The varieties of religious experience.} A study in human nature. London, New-York and Bombay. (1902). James is among the most outstanding thinkers of our time. He unites extensive knowledge with great power of observation, sharp criticism with idealistic enthusiasm, freedom from prejudice with earnest conviction. The most important of his writings is \textit{Principles of psychology} (1890), which on several points of psychology brings new points of view to bear and throughout discusses the questions in a spirited and stimulating way. In a collection of essays that has appeared under the title \textit{The will to believe} and other essays in popular philosophy (1897), he has in somewhat freer form expressed his peculiar union of philosophy of experience and a pronounced idealistic view of life, and already here he enters upon questions of the philosophy of religion. The new work is wholly devoted to the religious problem, in particular to the purely psychological illumination of it, yet in such a way that the other points of view from which the problem must be treated also come forward. James makes use of the specialized works of the younger American authors, but in addition draws on an extensive literature of a biographical, especially autobiographical kind. He sees clearly that the subject is too manifold to be exhausted in a single presentation, and therefore --- also in order to escape a dispute about words over what religion is --- he declares that for the purposes of this his investigation he will keep to a definite meaning of the word religion or to a definite side of religion, and he then defines religion as: the feelings, acts and experiences of individual human beings, in so far as these human beings in their solitude conceive themselves as standing in relation to the divine, \opage{70}however they may conceive this. This personal religion he sets in opposition to institutional religion with its more or less developed theology or mythology and its ecclesiastical organization. He maintains that \textit{institutional religion} (with all that belongs to it) is always derived in relation to \textit{personal} religion; that follows already from the fact that it builds on tradition. If one were to keep to institutional religion alone, the investigation would come to concern merely worship, sacrifice, ceremonial and ecclesiastical organization, and one could then define religion as the art of winning the favor of the gods. In personal religion, by contrast, it is man's own inner states, his conscience, his worth, his helplessness or his imperfection, on which the weight falls. If one keeps to this side of religion, the concept of religion acquires a greater scope than when one regards the institutional as the main thing.

Personal religion is the way in which man reacts to life as a whole (a man's total reaction upon life), in opposition to the momentary and special reactions. The experiences we have had of life and of the world act upon our individual temperament and make us either energetic or indifferent, either pious or mocking, either dejected or enthusiastic toward life as a whole. Our ``reaction,'' which often takes place involuntarily and half unconsciously, expresses itself in the most complete answer we can give to the question: What is the character of this existence to which we belong? --- Yet James will call such a reaction religious only when it is borne by a serious and exalted mood, a mood that at its height leads to the giving up of one's own wishes, to willing submission to the pain of reality, because one feels oneself carried by a power that is greater than the one over which one has command oneself. There is in all pronounced religious feeling a tragic element, which rests on the fact that a higher feeling of happiness has displaced a lower. The religious feeling contains or presupposes so many oppositions that it is one of the richest human feelings. The world is richer for there being a devil when there is a St. Michael who can set his foot on his neck. This distinguishing mark naturally appears most clearly in the extreme forms of religion, forms that \opage{71}to the ordinary, everyday conception stand as abnormal; but in all forms of religious consciousness it will occur to a certain degree.

This description indicates the general frame within which James moves in his investigation. For a closer characterization of his work I shall first show how the three points of view that --- as mentioned in the Introduction --- must come forward in every comprehensive discussion of the religious problem also appear in James, although he has mainly kept to the purely psychological mode of consideration. Next I shall show how he conceives the nature of religious experience, what types of religion he finds, and to what fundamental idea religion, in his conception, can be traced back.

\underafsnit{a. The Three Points of View.}{a. The Three Points of View}

James designates his standpoint, both in a general philosophical respect and in a philosophy-of-religion respect, as the standpoint of experience. He wants to be a thoroughgoing empiricist. All intellectual operations --- whether they consist in construction, in comparison or in criticism --- presuppose immediate experience. He allies himself with the old English school, the one that began with John Locke, and which always demanded a definite accounting for the experiences on which we build our concepts. This school ought in his opinion far sooner than Kant to have the credit for the critical method in philosophy, the only method that makes philosophy a study that can be pursued by serious human beings. But it is a matter not only of the experiences we have had, but also of those we must expect or consistently bring about ourselves, if our thoughts are to be valid. The value of our thoughts rests on the conclusions we can draw from them. If all the propositions we set up were practically indifferent, i.e. if nothing at all followed from them, could we then distinguish between true and false thoughts? In vain has dogmatism sought for criteria of truth that could dispense with appealing to the future. In opposition to dogmatism stands here pragmatism, which tests every proposition by pursuing its consequences. This conception (in which James allies himself with the \opage{72}American thinker CHARLES PEIRCE, from whom the name pragmatism also derives) recalls the biological-economic theory of knowledge that has been developed in Germany by Richard Avenarius and Ernst Mach.

James's empirical standpoint leads him to skepticism toward speculative and theological ways of viewing. He asserts in particular multiplicity in existence and is critically disposed toward monism in its various forms. Pluralism, the doctrine of multiplicity, will, he believes, come more and more into honor --- both in science and in religion. From experience one can reach nothing universal, no absolute and all-embracing unity. And different readings might hold for different regions of existence. Religiously regarded, it is here especially the experiences of evil that become of significance: has not the Book of Job once and for all demonstrated the impossibility of solving the problem to which these experiences lead? Here the monistic hypothesis will always present great difficulties, and would it not then be the simplest way out to assume that the world from the first has included elements that were not in harmony with the whole? Ordinary people have always been more or less polytheistic, and it is also clear that from no experience can one derive the infinity of God. What we can experience can at most be this, that we stand in inner connection with something that is greater than we ourselves; but that this is an infinite cannot be demonstrated.

What thinking can accomplish in the religious domain can only be this: to order, determine and interpret the facts, not to produce them. Philosophy is here always secondary. Yet it can for that reason well be of use. For the religious consciousness always expresses its experiences in forms that derive from the intellectual tradition under whose influence it stands. Therefore only critical comparison can find what is really immediate and essential and separate it from what is local and accidental. At the same time it is the task of philosophy to exclude assumptions that are inconsistent or decidedly in conflict with scientific experience. What survives this purgatory can then be regarded as hypotheses whose justification and significance can become the subject of closer discussion, once they have been reduced to their simplest forms.

\opage{73}So much for the consideration from the standpoint of the theory of knowledge. However energetically James asserts it, it is nevertheless the psychological mode of consideration that he is most interested in, and which gives his work its greatest significance.

The psychology of religion makes important contributions to general psychology. There are many psychological phenomena that can be studied particularly well in the domain of religious life. James emphasizes, for instance, that the religious phenomena demonstrate how small a part of our spiritual life can be explained clearly and distinctly. Consciousness here fades through many degrees into the unconscious, or, as James prefers to say, into the subconscious. Immediate impulses and presentiments stir, which often, without our noticing it ourselves, give us the first premises on which clearly conscious thoughts are grounded. Conscious arguments often move only on the surface of our being, and an involuntary and immediate conviction is the deepest thing in us (the deep thing in us). Further, the study of religious life teaches us to know great differences between human beings, both with respect to the kind of life of feeling and with respect to the degree of receptivity, --- differences that we should otherwise hardly become so aware of. We learn how development toward one and the same end can proceed very differently, now continuously, now by leaps. Already purely theoretically or from the standpoint of natural history a psychological study of religious life has its significance.

The psychological mode of consideration then in turn becomes significant for the understanding of the religious problem. The religious processes are seen as special forms of processes that are also known in other psychical domains, and obey laws that are also known elsewhere. James here points to Starbuck's investigations of conversion in its connection with organic and psychical development in the transitional years. Quite special weight is laid by James on what moves beneath the ``threshold'' of consciousness or outside its ``Margin.'' He finds it very difficult here to draw a definite boundary and rejects atomistic psychology, for which consciousness is composed of single elements, each fixed and definite in itself. In any case the subliminal or ultramarginal influences make themselves \opage{74}constantly felt, and James for his part is inclined to regard these influences as the paths by which a higher order of things works in our inner life. These inner manifestations constitute our nearest, our real world; so inwardly do we not belong to the outer world. Here James finds a point of unity between religion and psychology, although he declares every closer determination of it to be an interpretation that is no longer purely psychological. Interpretation is not the decisive thing either; it is always derived, and one and the same experience can be the object of different religious interpretations\textsuperscript{12}). What most people especially lack on this point is criticism and caution, not faith (faith). They are too inclined to let dogmatic assumption (belief) follow headlong upon every vivid idea, especially when this is supported by involuntary interest. ---

The third point of view, the ethical, comes forward in the appraisal of the religious phenomena. Just as definitely as James distinguishes between experience and interpretation, just as definitely he distinguishes between description or explanation of what is experienced and valuation of it. Even what, from a purely medical point of view, must be called abnormal can by its content have great value. The value of a state does not rest on how it came to be. When we call certain states higher than others, it is not on account of what we think we know about their organic conditions. If there is anything of value in S\textsuperscript{ta} Teresa's thoughts, then it is immaterial whether she was hysterical or not. We evaluate a force by its effect, not by its source.

In the valuation no speculative or theological system can be taken as a basis. Here too we must be empiricists. We judge the religious phenomena by their fruits, and that is what human beings have really always done. A worship ceases when it exerts no effect on the mind, and when in its whole character it comes into conflict with something whose value we have experienced so thoroughly that we cannot deny it. The gods to which human beings hold fast are those they can use, and whose commandments confirm the demands they make on themselves and on others. We always use human standards (human standards). We ask whether the religious life proves itself as an ideal \opage{75}form of human life. It is of course a standard that varies with the times, but another standard we simply do not have. Only the testing must be repeated again and again, since under new conditions it might lead to new results. And the valuation will often turn out differently for different individuals. When psychological investigation demonstrates that there are great differences between human beings in their innermost nature, one must also assume that they need different spiritual nourishment, so that they come to stand differently in religious respect.

Toward religion on the whole James stands decidedly sympathetic. The best fruits of religious experience are in his conviction the best things history has to show. He feels here the inner life stirring with seriousness and energy, working with an inwardness and a concentration, in opposition to all that can inhibit and divide. It is as if one were raised up into a higher and purer atmosphere. --- James refers on this point to Sainte Beuve's valuation of the religious personalities he has studied in his famous work on Port-Royal. ---

We thus find in James the three main points of view in philosophical discussion of religious questions. He announces that in a later work he will counter the attempts that even in our day have been made at a more speculative treatment of the philosophy of religion, --- attempts that, strangely enough, occur more frequently in England and America than on the Continent. For the time being his work stands as a significant contribution in favor of the view that it is precisely those three points of view that are to be applied in the philosophy of religion. --- Unfortunately, in my characterization of James's standpoint I cannot let the openness, freshness and linguistic art he commands come into their full right.

\underafsnit{b. The Nature of Religious Experience.}{b. The Nature of Religious Experience}

Religious feeling is according to James not any special or elementary feeling. Every feeling can under certain conditions acquire a religious character, namely when it co-operates in the reckoning of the experiences of human life that James calls ``the total reaction toward life.'' A somewhat closer determination of what religion is is given by the statement that religion is \opage{76}one of the ways by which unity can come into our life. With unity comes a feeling of peace and happiness that surpasses all that can otherwise be experienced of that kind. This unity can be reached by sudden breakthrough or through a long development; here the individual differences come to light. But the main point is that the individual, who before has felt divided and inhibited, low and unhappy, now has the feeling of harmony, freedom and elevation. Thoughts that before stood as ideal possibilities out on the periphery of the soul have now become the central possession of the soul. It is as if new energy had streamed into the individual from sources that lie outside its consciousness, and which it would not be able to release directly by conscious effort. This subordinate significance of the conscious and voluntary (the hopeless inferiority of voluntary to instinctive action) in relation to the unconscious and involuntary is a characteristic trait in all religiosity. Morbid forms arise when intellectual power and interest are lacking to shape and interpret such experiences.

Religious experience has a mystical character, partly because of ineffability, since each person here can understand the others only indirectly, partly because of immediacy, since discursive thinking plays no role, partly because of the short duration of the states and the difficulty of recalling them, partly finally because of the passive relation in which the individual stands to these experiences. The non-religious life of consciousness knows similar states. The religious character arises, as already remarked, from the fact that man's relation to life as a whole is settled in such states; thereby they acquire a unity and energy that is otherwise lacking. What is specifically called mysticism is only a particular and often extreme form of what appears in all religious experience.

\underafsnit{c. The Most Important Types.}{c. The Most Important Types}

James declares that the most important result he has reached in his investigations of religious consciousness is the great differences that it presents. A significant part \opage{77}of his work therefore naturally consists in tracing these differences back to certain types.

There is in particular one opposition that he emphasizes very strongly, namely the one between the healthy and the sick mind, or --- to use an expression he has from Francis Newman --- between the once born and the twice born (the once-born and the twice-born). The sick souls need a change, a crisis, a conversion in order to reach unity and peace, whereas there are healthy souls who by immediate unfolding and surrender reach the goal. The religion of the healthy James designates as naturalism, that of the sick as salvationism. But within each of these groups very different personalities and tendencies can in turn appear.

The religiosity of the healthy souls is an enthusiasm that holds fast to the bright side of existence and perhaps even declares the dark sides to be illusions. Even if they do not lack an eye for what inhibits and darkens, they are nevertheless assured that it can be overcome; so great is their immediate trust in inner energy and in the harmony that prevails in the universe. This tendency appears more in Greek antiquity than in Christianity, more in Catholicism than in Protestantism, more among the Latin races than among the Germanic. In more recent times it appears in Spinoza and Rousseau, in the liberal tendency within Christianity, especially in Unitarianism (Emerson, Parker). In America it stands in sharp opposition to Methodism and the Revival movement. In the most recent time the so-called mind-cure movement is a characteristic example. It heals both soul and body by constantly instilling and holding fast to the thought that pain is an illusion, and that everything in reality is bright and glorious. Fear and worry are banished. It works through suggestions, not least auto-suggestions. By constantly shouting Hurrah for the universe one strengthens oneself in the belief that the universe deserves a Hurrah.

The sick soul succumbs to pain and disharmony. The threshold of consciousness toward that kind of experiences is once and for all lowered, whereas in the healthy soul it is raised. The world stands as empty and vain; fear is nourished of \opage{78}what the course of life may bring; the soul feels itself sinful and corrupt. The feeling of emptiness, the feeling of fear and the feeling of sin are the three main forms of mental illness. In Tolstoy the feeling of emptiness prevails, in Bunyan the feeling of fear and the feeling of sin. There is here an original disharmony in the temperament that demands a second birth. And even if this is reached, natures of this kind will hardly attain the overflowing feeling of life that the healthy can have.

There exist of course in reality a multitude of transitional forms between the two types. But in their extreme forms the opposition between them is so great that they hardly understand each other. Each will regard its view of life as the one right one. Thus Wesley maintained conversion as the only way to peace, whereas Emerson ascribed to the nature that needed only one birth a higher value than to the one that had to struggle hard with the devils in its inner life. James himself thinks that the experiences of life on which the twice born can build are deeper and more complete than those of the first type. No prophet can bring a decisive message to human beings if he does not say things that can ring as realities for men like Bunyan and Tolstoy. The most complete religions must be those that both know and overcome the darkest elements of life. Therein rests the superiority of Buddhism and Christianity. ---

In relation to this great opposition the other oppositions that can be pointed out are for James of less importance. Thus the opposition between leaping and continuous development, --- the opposition between natures in whom the elements that are to be overcome must first have lost their vital force, and other natures in whom the new life can break through before the old is wholly exhausted, --- and the opposition between a development that demands conscious effort (even if this does not lead directly to the goal) and a development that runs its course throughout involuntarily. ---

Occasionally James gives, besides this general doctrine of types, also hints of his own purely individual faith. He has, in a state of nervous slackness that long prevailed, sought help from the mind-cure movement and believes himself to have been cured \opage{79}by its help, which does not prevent him from taking a critical stance toward it. When he states that mysticism gives only in extreme form what all religious experience contains, he adds that he himself does not personally know mystical states. His faith he characterizes by saying that although he cannot join himself to the current Christianity or to speculative theism, he is convinced that in those states in which man feels himself in the most inward unity with the Highest he knows, new power comes into the world and new points of departure are set; so far is his faith supernaturalistic, assuming divine influences within the domain of natural experience. He adds that he knows very well that the prevailing opinion in academic circles runs in the opposite direction. --- Belief in immortality plays only a subordinate role for him (although in a special little work he has sought to show that it cannot be refuted). He says of this: If only ``in all eternity'' provision is made for our ideals, I do not see why we should not be willing to hand over the care of them to other hands than ours.

\underafsnit{d. The Fundamental Religious Idea.}{d. The Fundamental Religious Idea}

Although religion according to James has its ground in the life of feeling and its experiences, and although it appears in very different forms, there is nevertheless a fundamental idea or a hypothesis on which it always rests. It is this, that a state of harmony, unity and peace can be reached, even if it must happen through struggles and crises, and that this end is reached by an influx of energy from a more comprehensive order of things, which psychologically regarded expresses itself in the unconscious nature in us. How we in particular interpret what happens here is not the essential thing; here purely individual supplements and elaborations (overbelief) will always come in. The main point is that for the religious consciousness tragedy, however often it may announce itself on the path of life, is always only provisional and partial. Shipwreck and dissolution are not the last word of existence. All religion thus rests ultimately on a cosmic hypothesis. On the other hand the persistence of religion is not dependent on any special dogmatic assumption whatever. Of course --- if religion is made one with or wholly dependent \opage{80}on an animistic conception of the world, then its time will be over, when animism dies out; likewise if one makes it into magic or fetishism. For James what is waged here easily becomes a dispute about words. The decisive thing for him is that the inner experiences on which religion essentially builds are the most concrete and immediate experiences we have. In relation to them everything we can experience about the outer relations of the world has an abstract and general character. There is no reason to assume that the time of such experiences should be over, even if they must be separated from connection with elements with which they have hitherto been inclined to fuse.

\underafsnit{e. Remarks on James’s Philosophy of Religion.}{e. Remarks on James’s Philosophy of Religion}

The study of James's book, which appeared a year after my own attempt in the philosophy of religion, has been of great interest to me. Not only is it long since I read a book that has so taught me to look with new and fresh eyes at life and at human beings as this book by the American philosopher has managed to do. But the book acquired a special significance for me in that the whole direction and way of treatment, despite all difference on not unessential points (e.g. in what James calls his ``supernaturalism''), is so closely akin. I have therefore a personal reason for calling attention to it. By its spirited form and its many interesting details it will moreover reward study, even if one stands further from the book's direction than I do\textsuperscript{13}). Only a few remarks will I put forward here in order to mark my standpoint toward it.

It is in my eyes with full right that James distinguishes between personal and institutional religion. It is admittedly a difference that is easier to make for the higher than for the lower religions, and it is perhaps for that reason that historians of religion sometimes shrink from admitting it. But psychological interest must, by the nature of the case, concern above all the personal element in religion, whereas the objective or ``institutional'' side has significance only as symptom. And what is the main point: the religious problem really concerns \opage{81}only the subjective side. The question of the position of religion --- now and in the future --- in human intellectual life is after all primarily whether man's feeling and need will be able to go on holding fast to the forms and symbols hitherto given, and whether, if these should prove insufficient, new ones can be found or created. He who fixes his gaze on the continuity in spiritual development must decidedly lay the main weight on the subjective, personal domain, however much this may be veiled by the historical predominance of the outer forms.

Yet it seems to me that James goes too lightly over the difficulties the problem of continuity here presents. It is by no means self-evident that the personal element in religion can be preserved if the historical forms fall away. As it will be seen from my Philosophy of Religion, I stand here in such a way that I see a \textit{possibility} of a preservation of the inwardness and concentration in the conception of life that has had so essential a support in the historical forms of religion, although I cannot adduce any compelling proof of the necessity thereof, --- not even that equivalents for it can be given. We face the possibility that value will here be lost. History shows us a constant reciprocal relation between the personal and the institutional. If this reciprocal relation is to persist in the future, it is possible only if the human spirit in the course of development retains the ability to find or create great symbols for the deepest experiences of life, symbols that can acquire common significance for larger or smaller groups of human beings. It can be hoped, but not proved. And yet it is one of the most important questions with regard to the future of our race. One must look at the greatest conflicts of life with a very light mind and a very unclear eye if one does not feel the sting of this problem. --- Whether under the new conditions one will be able to go on using the \textit{word} religion is an indifferent question, which James treats with the irony it deserves. ---

If James's description of religious experience is compared with the one I have given, there appears an agreement that shows that we have the same phenomena in view. Religious experience is for him characterized by being man's ``total \opage{82}reaction'' toward life; it is an experience of peace, unity and energy, often won through a crisis that led to the displacement of darkness and disharmony. What I miss here is a definite emphasis on the fact that the disharmony that religious consciousness in all its forms strives to dissolve is a disharmony between value and reality. By the relation between the values that man knows and the reality that he knows his religion is determined. What is valuable can for him be the physical life itself in its persistence and unfolding, or it can be ideal goods. His total feeling toward life is always determined by whether the heaven of values is darkened or not. It is from this point of view that I have come to my hypothesis that the persistence of value is the fundamental religious idea or the religious axiom. James lays the weight on the inner feeling of unity and on the experience of an influx of inner energy, without taking sufficient account of the fact that this unity and heightening of energy for man means that the disharmony between value and reality is removed. Yet I do not think that James will disagree with me in my hypothesis. As already mentioned above, he indeed suggests that religion, while built first and foremost on subjective experience, rests more or less consciously on the presupposition that the tragic in the world can only be partial and transitory, --- that shipwreck and dissolution cannot be the last word of existence. The religious consciousness thus comes by its nature to presuppose something about cosmic relations, about the innermost nature of existence. It is on this point that religion has always had a tendency to become metaphysics, just as conversely metaphysics, without always being conscious of it, has proceeded from religious presuppositions in its constructions. According to James every attempt at a more definite and concrete formulation of those ultimate presuppositions is ``over-belief,'' i.e. leads from the domain of thought and experience over into that of fantasy and mysticism. James's own ``supernaturalism'' is evidently an example of this. About the new beginnings he assumes, he has no right to say more than that their cause cannot for the time being be specified; that they \opage{83}should have a supernatural cause is not motivated according to his own conception. ---

There is still one essential point on which I agree with James, namely the great weight he lays on the individual differences with respect to religious experience and faith. His masterly portrayal of the opposition between the healthy and the sick souls' way of taking things is a great contribution in favor of what I call the principle of personality. This opposition corresponds to the opposition already mentioned above between expansive and disharmonious natures. But James has carried the characterization through with such force and clarity that the question whether different personalities can really have the ``same'' religion arises with still greater necessity than before. ---

\kapitel{}{Notes.}{Notes}\opage{84}
\streg

1) P. 1. --- Cf. my essay \textit{Filosofien og Livet.} (Tilskueren 1901).

2) P. 4. --- L. P. \textsc{Larsen}: \textit{Høffdings Filosofi blandt Indiens Studenter.} (Studenterhjemmets Julebog 1901). --- I cannot refrain from remarking that the pompous title this essay has been given awakens false ideas both of ``Høffding's Philosophy'', of ``India's Students'' and --- of the content of the essay itself. On the other hand it suits the character of the essay rather well that it has found its place in an edifying publication. The author says of me that my interest is more philosophico-religious than religious; this I can only regard as a compliment, since it was precisely a philosophy of religion that I meant to write. In return I may well say that his interest in the edifying is greater than his interest in the philosophical. Perhaps that is why he does not read philosophical books with the necessary attention either. Thus he says that on p. 214 of my Philosophy of Religion I ``protest against the dogma of creation''. No such protest is to be found in my book. At the place cited I discuss the difficulty that a great religious genius like Augustine found in the creation of a finite world, because a diminution of value would thereby take place, and I see in it a testimony that faith in the persistence of value is the fundamental religious idea. If the author of the essay does not acknowledge that difficulty, he must settle the matter with Augustine. For the present I hold to the old Father of the Church, who could both \textit{experience and think.} --- When I have here had to reject a theological critic's understanding of my book, it is all the greater a pleasure to me to refer to two theological critics, a Danish and a German, who both give a clear and exact account of the content of the work they set out to judge, just as their critical remarks too have been instructive to me on several points: F. C. \textsc{Krarup} in Dansk Tidsskrift, April 1901, and \textsc{Friedrich Lipsius} in Protestantische Monatshefte. 6. Jahrg. Heft 7.

3) P. 6. --- In the first edition of his \textit{System der Philosophie} (p. 444) \textsc{Wilhelm Wundt} had declared that philosophy could prove the necessity of faith (although it could not convert faith into knowledge). In the second edition this sentence is omitted, without any reason being given for it. Possibly the omission was occasioned by a growing awareness of the purely empirical mode of consideration at which we must here stop.

\opage{85}4) P. 10. --- \textsc{Dowden}: \textit{Puritan and Anglican.} London 1900. p. 234. --- Cf. my \textit{Religionsfilosofi.} p. 278.

5) P. 15. --- For the sake of clarity I set down Siebeck's statement in literal translation: ``When the personality, freely and consciously, rejects the traditional or involuntarily developed content of faith, declaring faith outright to be something that ought not to be, and holding the scientific procedure to be the only justified one, an illusion underlies this. The personality does not become aware that this act, which is an act of freedom, becomes possible only through a moment that belongs to the essence of faith .... For faith in a highest good, and in a value that can through it fall to the share of the personality, always lies beforehand at the basis of the decision; only, the highest good is in this case placed exclusively in the persistence and coherence of scientific knowledge. But this last assumption is not itself first the result of a scientific proof; for without it there would never have arisen any impulse to scientific activity''. H. \textsc{Siebeck}: \textit{Religionsphilosophie.} p. 200.

6) P. 19. --- Cf. with this \textsc{Schleiermacher's} conception of prayer \textit{(Der christliche Glaube} § 146--147). A theory of prayer that presupposes a reciprocal relation between the Creator and the created denotes, according to Schleiermacher, a transition to the magical. There is an inner connection between prayer and the answering of prayer; both are at bottom one, and likewise in true prayer passivity and activity are united. The main thing is that prayer itself belongs to ``the divine order of the world'' just as much as its being answered or not answered. (ibid. § 47). --- An American physician has recently expressed himself thus about prayer: ``Perhaps the very act of prayer is the way by which I am helped by the divine power: when I speak to God, it is God who speaks in me. It is strange that, although I do not believe in miracles, either in the past or in the present, I nevertheless still believe that God must have some way or other of realizing my ideal of divine humanity''. (Quoted by \textsc{Leuba}: \textit{The content of religious consciousness.} The Monist. July 1901. p. 553).

7) P. 21. --- In a certain sense all scientific concepts too are in the end symbolic. On the difference between scientific symbolism and religious (or metaphysical) symbolism see \textit{Religionsfilosofi} p. 70--72. --- \textit{Filosofiske Problemer} (University programme 1902). p. 37 f.

8) P. 37. --- Cf. \textit{Det psykologiske Grundlag for logiske Domme.} p. 47 f. --- \textit{Filosofiske Problemer.} p. 75 f.

9) P. 38. --- Cf. with this the younger \textsc{Ritschl's} work \textit{Ueber Werturteile} (Freiburg 1895). Chap. 5 (Die geschichtlich gegebenen geistigen Werte im Unterschied und im Verhältniss zu der subjectiven Wertbeurteilung).

10) P. 40. --- See \textit{Religionsfilosofi.} p. 111--115; 251--257. Cf. also \textsc{Leuba's} criticism of Murisier in Revue philosophique. Juillet 1902 p. 33--35.

11) P. 54. --- Perhaps a distinction might here be drawn between the leaders of the Church and its believers. Its leaders have often been decidedly choleric and active natures; but the believers as such were referred to a predominant passivity, and it was the moods of sorrow and of hope that were cultivated in them through the cult.

12) P. 74. --- On this point James joins Leuba. ``Leuba is undoubtedly right in contending that the conceptual \textit{belief} . . . although so often efficacious and antecedent, is really accessory and non-essential, and that the joyous conviction can also come by far other channels than this conception. It is to the joyous conviction itself, the assurance that all is well with one, that \opage{86}he would give the name of \textit{faith} par excellence.'' \textit{(Varieties of religious belief.} p. 246 f.). --- Cf. with this \textit{The Will to believe.} p. X: ``What mankind at large most lacks is criticism and caution, not faith. Its cardinal weakness is to let belief follow recklessly upon lively conception, especially when the conception has instinctive liking at its back''. James adds here that if he were speaking to the Salvation Army or similar circles, he would not stress the will and the right to believe; such circles would rather need ``the fresh north-west wind of science'' to sweep their morbidness and barbarism away. But he is speaking to academic circles, whose needs are different: ``Paralysis of their native capacity for faith and timorous abulia in the religious field are their special forms of mental weakness''.

13) P. 80. --- Of philosophical criticisms of James's book I know only two, one by the Geneva psychologist \textsc{Th. Flournoy} in Revue philosophique (Nov. 1902) and one by J. H. \textsc{Muirhead} (Professor in Birmingham) in International Journal of Ethics (January 1903). Flournoy joins James with enthusiasm; Muirhead acknowledges him as a psychologist, but is dissatisfied on behalf of metaphysics.

\end{document}
